\documentclass[11pt,twoside]{book}

\usepackage[T1]{fontenc}
\usepackage[utf8]{inputenc}
\usepackage{mathpazo}
\usepackage{courier}
\usepackage{microtype}
\usepackage[margin=1.25in]{geometry}
\usepackage{amsmath,amssymb,amsthm}
\usepackage{xcolor}
\usepackage{listings}
\usepackage{booktabs}
\usepackage[hidelinks]{hyperref}

\definecolor{codebg}{gray}{0.96}
\definecolor{codecomment}{rgb}{0.30,0.42,0.30}
\definecolor{codekeyword}{rgb}{0.10,0.20,0.55}
\definecolor{codestring}{rgb}{0.55,0.22,0.12}

\lstdefinelanguage{GoLang}{
  morekeywords={break,case,chan,const,continue,default,defer,else,fallthrough,for,func,go,goto,if,import,interface,map,package,range,return,select,struct,switch,type,var},
  morekeywords=[2]{bool,byte,error,float64,int,int64,string,rune,uint8,any,true,false,nil,iota},
  morecomment=[l]{//},
  morecomment=[s]{/*}{*/},
  morestring=[b]",
  morestring=[b]`,
  sensitive=true
}

\lstdefinestyle{gostyle}{
  language=GoLang,
  basicstyle=\ttfamily\small,
  keywordstyle=\color{codekeyword}\bfseries,
  keywordstyle=[2]\color{codekeyword},
  commentstyle=\color{codecomment}\itshape,
  stringstyle=\color{codestring},
  backgroundcolor=\color{codebg},
  frame=single,
  rulecolor=\color{black!25},
  numbers=left,
  numberstyle=\tiny\color{black!50},
  numbersep=8pt,
  tabsize=4,
  showstringspaces=false,
  breaklines=true,
  breakatwhitespace=false,
  columns=fullflexible,
  keepspaces=true,
  captionpos=b,
  xleftmargin=1.6em,
  framexleftmargin=1.2em,
  aboveskip=1.2em,
  belowskip=0.8em
}

\lstdefinestyle{shellstyle}{
  basicstyle=\ttfamily\small,
  backgroundcolor=\color{codebg},
  frame=single,
  rulecolor=\color{black!25},
  showstringspaces=false,
  breaklines=true,
  columns=fullflexible,
  keepspaces=true,
  captionpos=b,
  aboveskip=1.2em,
  belowskip=0.8em
}

% Source file listing: path relative to the golab module, caption, label.
\newcommand{\golisting}[3]{%
  \lstinputlisting[style=gostyle,caption={#2},label={#3}]{../golab/#1}}

\theoremstyle{plain}
\newtheorem{proposition}{Proposition}[chapter]
\theoremstyle{definition}
\newtheorem{definition}{Definition}[chapter]
\newtheorem{exercise}{Exercise}[chapter]
\theoremstyle{remark}
\newtheorem*{remark}{Remark}

\newcommand{\golab}{\texttt{golab}}
\newcommand{\gocode}[1]{\texttt{#1}}

\title{\textbf{Go as an Instrument of Continuation}\\[0.6em]
\large From Local Workflows to Concurrent Systems, Containers, and Governed Services}
\author{Flyxion\\Independent Researcher}
\date{October 2026}

\setlength{\emergencystretch}{2.5em}

\begin{document}

\frontmatter
\maketitle
\tableofcontents

\chapter{Preface}

This book teaches the Go programming language by building one system from the first line to the last, instead of presenting a sequence of unrelated exercises. The system is a laboratory tool called \golab{}. In its first form it is a command-line program that finds files in a research directory, reads them and records their cryptographic digests, runs external programs on them, and writes each of these actions into a ledger whose entries are chained together so that later alteration can be detected. Later parts of the book give the same tool a concurrent worker pool, durable indexes, an HTTP interface, container packaging, and a set of experimental harnesses. At every stage the program that exists is the program the chapter describes, and it compiles, passes its tests, and runs.

The conventional route into Go begins with variables and loops and arrives at web services by way of a long tour of the standard library. The route taken here is different because the intended reader is someone who already works fluently at a shell, who has written Bash scripts that traverse directories and invoke converters, and who has felt the point at which such scripts stop being adequate. The point is rarely a matter of speed. It arrives when a workflow must keep distinctions that the shell leaves implicit: which files were merely found and which were actually read, which outputs are evidence and which are interpretations of evidence, which failures left the filesystem partly changed, and which actions were permitted rather than merely performed. Go is a good language in which to write such distinctions down, because its types are simple enough to be read quickly, its errors are ordinary values that must be handled in plain sight, and its tools produce a single executable that can be run from the same shell the reader already uses.

Several ideas that organize the author's wider research recur here as questions that shape the design of \golab{}: the difference between a proposed transition and an accepted one, the difference between recording an event and authorizing it, the difference between a source and a summary of that source, and the difference between routing work to a destination and adjudicating whether the work should be done. These ideas are used as structural analogies and as prompts for design. The language does not enforce any of them automatically, and the book takes care to say, at each point, whether a given property is guaranteed by the compiler, guaranteed by a library under stated assumptions, or merely maintained by convention and discipline. A reader who finds that a chapter makes a stronger claim than its code supports should treat that as a defect of the chapter.

The book is organized in six parts. Part~I, Local Computation, runs through six chapters. It introduces programs as transition systems, values and types, filesystem traversal and hashing, errors as values, the execution of external programs, and the hash-chained ledger. Part~II, Governed Workflows, runs through three further chapters and introduces interfaces and composition, together with a five-stage record pipeline in which storing, validating, admitting, projecting, and activating a record are separate steps with separate types. Part~III, Concurrent Execution, runs through three more chapters and gives the tool a bounded worker pool, cancellation through contexts, and a ledger writer that serializes appends. Part~IV, Memory and Reconstruction, runs through three more chapters and adds a file that is replaced atomically and an index that is derived from the ledger and can always be rebuilt. Part~V, Distributed Systems, runs through three more chapters and gives the tool an HTTP interface, separates routing from adjudication and from admission, and runs the service in a minimal container. Part~VI, Experimental Systems, runs through four chapters and treats the tool as an object of study, with benchmarks and profiles, structured logs and injected faults, manifests that fix what an experiment ran on, and a search for what a summary of a ledger discards. All six parts are complete in this volume.

Every listing in Parts~I to VI is included directly from the source files of the \golab{} module, so the text and the code cannot drift apart. The module was built and tested with Go 1.24 on Linux, and nothing in the whole book depends on a feature newer than Go 1.21. Readers working under the Windows Subsystem for Linux can follow the book without modification. Commands that appear in a shell listing were run, and the output shown is the output that was produced. Appendix~\ref{app:guarantees} collects, in one table, which properties the compiler or toolchain guarantees and which are conventions, so that no claim in the text need be read as stronger than that table allows.


\mainmatter

\part{Local Computation}
\chapter{A Program Is a Controlled Transition}
\label{ch:transition}

\section{The Case for a Compiled Tool}

Go occupies a particular position among programming languages. It compiles to native executables, it manages memory automatically, it has structural interfaces, explicit error handling, and lightweight concurrency, and it ships with a standard library broad enough that many useful programs need no outside dependencies at all. Its syntax is deliberately restrained, which means that a reader can learn most of the language in a short time, but its operational reach is large. A program written in Go can traverse a filesystem, launch and supervise other programs, transform data, serve network requests, and run the same way on a laptop, in a scheduled job, and inside a container.

For a reader accustomed to Bash, Go at first looks unnecessarily elaborate. A shell script can list files, filter text, and invoke a converter in a handful of lines, whereas the equivalent Go program needs declarations, a compilation step, and explicit handling of every error. The elaboration pays for itself when a workflow has to preserve distinctions that a shell leaves unstated. Consider a research directory that holds source documents, generated PDFs, transcripts, measurements, and intermediate files. A script can find the files and call a converter on each. Suppose, however, that one conversion fails, that an output is older than its source, or that a machine-generated summary has been mistaken for the document it summarizes. The difficulty is no longer that of executing commands. It is the difficulty of maintaining a relation between operations, evidence, and the actions that the evidence permits afterward, and a tool that is expected to maintain such a relation needs a way to name its states, to refuse transitions that should not occur, and to leave a record of what it did. The shell remains the better instrument for brief interactive composition of commands, and nothing in this book proposes that it be abandoned. The proposal is narrower: when a sequence of convenient commands has become a system whose intermediate states and failure modes matter, it is time to write the system down in a language that can hold those states in its types.

\section{The Smallest Program}

The first program prints one line. It is reproduced here exactly as it exists in the \golab{} module, under \gocode{examples/ch01\_hello}.

\golisting{examples/ch01_hello/main.go}{The minimal program.}{lst:hello}

A Go source file begins by naming its package. A package named \gocode{main} that defines a function named \gocode{main} is a program, and the function is where execution begins. The \gocode{import} declaration makes the standard package \gocode{fmt} available, and \gocode{fmt.Println} writes its arguments to standard output followed by a newline. Executed with \gocode{go run}, the program compiles and runs in one step.

\begin{lstlisting}[style=shellstyle,caption={Running the minimal program.}]
$ go run ./examples/ch01_hello
Continuation begins.
\end{lstlisting}

Unlike a Bash script, which is read and interpreted by a shell each time it runs, a Go program is compiled to machine code before it runs. The \gocode{go build} command performs the compilation alone and leaves an executable behind, which can then be run from a shell, scheduled by another tool, or copied into a container image. For programs that depend only on the standard library, that executable is ordinarily self-contained, although its portability is still bounded by the operating system and processor architecture for which it was built. A module, which is the unit in which Go code is versioned and built, is created with \gocode{go mod init}; the laboratory tool of this book is the module named \gocode{golab}, and every path in the listings is relative to its root.

\section{Computation as a Transition System}

A program can be described abstractly as a transition. Let $S$ be the set of states the program can occupy, $I$ the set of inputs it can receive, $O$ the set of outputs it can produce, and $E$ the set of exceptional outcomes it can report. A single step of the program is then a map
\[
T : S \times I \longrightarrow S \times O \times E,
\]
which takes a state and an input and yields a successor state together with an output and a possible error. The description is not peculiar to Go, and it is deliberately coarse, since it says nothing about how the step is carried out. Its use is that it separates three things that a casual description of a program runs together: what the program remembers between steps, what it is asked to do, and what happens as a consequence.

The simplest nontrivial case is a counter. Its state is a single integer, and the operation it supports adds one to it. The version below is the example from the module under \gocode{examples/ch01\_counter}.

\golisting{examples/ch01_counter/main.go}{A counter whose transition returns a new state.}{lst:counter}

The program prints \gocode{0} and then \gocode{1}. The point of the example is the manner in which \gocode{Increment} treats its argument. In Go, an argument of struct type is passed by value, so the function receives a private copy of \gocode{State}, increments the copy, and returns it. The variable \gocode{current} in \gocode{main} is not changed. The step therefore has the shape $s_{t+1} = T(s_t)$ with the old state still available for comparison, which is exactly the shape of the abstract definition. A program written in this style makes each transition visible as a function call whose input and output can both be inspected, and that visibility is what later chapters rely on when they attach a policy to a transition or a record to its outcome.

\begin{remark}
Process-oriented accounts of computation, including the SpherePOP calculus that appears elsewhere in the author's work, ask of any system not only which state it is in but which transitions it permits, which consequences of a transition persist, and which earlier histories remain reachable. That vocabulary suggests the questions this book asks of \golab{}. It does not follow that Go implements any such calculus. Go supplies functions, values, and types, and the distinctions of a richer theory have to be encoded by the programmer in those materials.
\end{remark}

\section{Admissibility}

Not every proposed transition should be accepted. Suppose the counter must never become negative. The requirement can be stated as an admissibility condition, a function $A$ that inspects the current state and the proposed input and answers whether the transition may proceed:
\[
A(s,i) =
\begin{cases}
1 & \text{if the proposed transition is permitted},\\
0 & \text{otherwise.}
\end{cases}
\]
The transition is then defined only where $A(s,i)=1$. Where $A(s,i)=0$ the system reports a refusal and the state is left exactly as it was.

Go expresses refusal through ordinary return values. A function that can fail returns its result together with a value of the built-in type \gocode{error}, and the value \gocode{nil} means that no error occurred. The example below, from \gocode{examples/ch01\_decrement}, applies the convention to the counter.

\golisting{examples/ch01_decrement/main.go}{A decrement that refuses to produce a negative state.}{lst:decrement}

Run with a counter of zero, the program prints \gocode{Rejected: negative state prohibited}. Two features of the example deserve attention. The first is that the distinction between a proposed transition and an accepted one is visible in the control flow: the call to \gocode{Decrement} proposes, the test of \gocode{err} decides, and only the branch in which \gocode{err} is \gocode{nil} proceeds to use the new state. The second is that on refusal the function returns the original state unchanged, so a caller who inspects the returned state after a failure sees the state from before the attempt. Both features are conventions of Go programming and neither is enforced by the compiler. A caller is free to write \gocode{next, \_ := Decrement(current)} and continue as though nothing had been refused, and the compiler will accept the program. Static analysis tools outside the compiler can flag discarded errors, but the language itself does not require that an error be examined. This gap between reporting a failed transition and enforcing an admissibility boundary is the subject of the fourth exercise at the end of this chapter, and it returns in Chapter~\ref{ch:errors}.

A more elaborate system will distinguish validation, authorization, and commitment, which are three different questions that are easily collapsed into one function. Validation asks whether the proposed input is well formed. Authorization asks whether the actor is entitled to request it. Commitment asks whether the resulting change is made durable. Part~II returns to these separations. For now it is enough to note that the single function \gocode{Decrement} answers only the first of them and that nothing in the signature says so.

\section{Admissibility Without Overflow}

A bounded counter exposes a hazard that every program with arithmetic must face. Fix a maximum value $M$ and require that the counter stay in the interval $[0, M]$. Adding a caller-supplied amount $\delta$ is then admissible when the sum lies in that interval. The obvious test computes $c+\delta$ and compares the result with the bounds, but in Go the signed integer types wrap around on overflow without any error or panic, so a sum that exceeds the largest representable integer silently becomes a large negative number, and a test that was meant to protect the bound can be defeated by the very arithmetic it performs. The test must therefore be arranged so that it never computes a quantity that can overflow.

\golisting{internal/transition/transition.go}{The transition package, with an overflow-safe admissibility test.}{lst:transition}

The function \gocode{Add} in Listing~\ref{lst:transition} applies the principle on both sides of the interval. When $\delta \ge 0$ the transition is admissible exactly when $\delta \le M - c$, and the right-hand side lies between $0$ and $M$ for every legitimate state $c$, so it cannot overflow. When $\delta < 0$ the transition is admissible exactly when $\delta \ge -c$, and since $c \ge 0$ the quantity $-c$ is always representable. The following proposition records the argument.

\begin{proposition}\label{prop:add}
Let $0 \le c \le M$ and let $\delta$ be any integer in the representable range. Then $0 \le c+\delta \le M$ if and only if either $\delta \ge 0$ and $\delta \le M-c$, or $\delta < 0$ and $\delta \ge -c$. Moreover, evaluating $M-c$ and $-c$ cannot overflow.
\end{proposition}

\begin{proof}
If $\delta \ge 0$ then $c+\delta \ge 0$ automatically, and $c+\delta \le M$ holds exactly when $\delta \le M-c$. If $\delta < 0$ then $c+\delta \le M$ automatically, and $c+\delta \ge 0$ holds exactly when $\delta \ge -c$. Because $0 \le c \le M$, the value $M-c$ lies in $[0,M]$ and the value $-c$ lies in $[-M,0]$, and both intervals are contained in the representable range whenever $M$ is.
\end{proof}

The reason for writing the lower-bound test as $\delta < -c$ and not as $-\delta > c$ is not cosmetic. In two's complement arithmetic the most negative integer has no positive counterpart, so negating it yields itself, and the apparently equivalent test accepts exactly the input that it should reject most firmly. The first version of \gocode{Add} written for this chapter used the negated form, and it was the table-driven test in Listing~\ref{lst:transition-test}, which includes the most negative integer as a case, that rejected it. The episode is retained in the text because it illustrates the role of tests in the method of the book: a test is the means by which a claim about admissibility is checked against the arithmetic of the machine on which it must hold.

\golisting{internal/transition/transition_test.go}{Table-driven tests for the transition package.}{lst:transition-test}

The test file shows the standard shape of a Go test. A function whose name begins with \gocode{Test} and which takes a pointer to \gocode{testing.T} is run by the \gocode{go test} command. The table in \gocode{TestAdd} lists cases as data, and each case is run as a named subtest, so that a failure identifies the case that failed. The comparison \gocode{errors.Is(err, c.err)} asks whether the returned error is, or wraps, the expected sentinel value, a technique taken up in Chapter~\ref{ch:errors}.

\section{From Transitions to Workflows}

The step from a counter to a workflow is a change of subject matter and not of method. Suppose that the task is to list the regular files below a directory. In Bash the answer is a single command, \gocode{find . -type f}. The corresponding Go program, from \gocode{examples/ch01\_walk}, is considerably longer.

\golisting{examples/ch01_walk/main.go}{Listing regular files below the current directory.}{lst:walk}

The program hands \gocode{filepath.WalkDir} a function to be called for every entry. The function receives the path, a description of the entry, and any error encountered in reaching it, and it returns an error to stop the traversal or \gocode{nil} to continue. The length of the program relative to the shell command is real, but what the extra length buys is a place to attach things the shell command has no room for: a typed description of each file, a classification, a content hash, a policy about unreadable directories, and a record of the traversal. Chapter~\ref{ch:traversal} develops those additions, and by the end of Part~I the same traversal logic sits behind commands such as \gocode{golab discover}, \gocode{golab verify}, and \gocode{golab run}.

A tool of this kind has four responsibilities which can coexist in one executable but should remain distinguishable in its design. A command is an instruction that the tool accepts from its operator. An operation is the work performed in response to a command. A policy is the rule that determines whether the work is permitted. A record is the durable trace that preserves what happened. In the version of \golab{} developed in Part~I the commands are implemented in the \gocode{cmd/golab} package, the operations in the packages \gocode{artifact} and \gocode{runner}, the records in the package \gocode{ledger}, and the policy, which is still rudimentary, in the classification rules and in the admissibility checks of the transition package. Part~II gives policy a package of its own. The ordering of the book follows from this principle: distribution changes the circumstances in which computation takes place, but it does not remove the need to say what the computation is. Goroutines, channels, network services, and containers are therefore deferred until the local tool is understood, since each of them multiplies the consequences of whatever ambiguity the local tool still contains.

\section*{Exercises}

\begin{exercise}
Modify the counter of Listing~\ref{lst:counter} so that the caller supplies an increment amount, and reject any transition that would exceed a fixed maximum. Write the admissibility test so that it cannot overflow, state the argument for its correctness in the manner of Proposition~\ref{prop:add}, and add table-driven tests that include the largest and smallest representable integers.
\end{exercise}

\begin{exercise}
Modify the traversal program of Listing~\ref{lst:walk} so that it prints only files whose names end in \gocode{.go}, and then add a count of the matching files that is reported after the traversal completes. Decide what the program should do when a subdirectory cannot be read, and justify the choice.
\end{exercise}

\begin{exercise}
Define two distinct Go types, one for a file that has been discovered and one for a file whose content has been verified, and consider which functions ought to accept each. Explain what, if anything, prevents a caller from manufacturing a value of the second type without performing any verification. Compare your answer with the treatment in Chapter~\ref{ch:types}.
\end{exercise}

\begin{exercise}
Explain why returning an error from a function does not by itself guarantee that the caller will respond correctly. State the difference between reporting a failed transition and enforcing an admissibility boundary, and describe one mechanism, inside or outside the language, that moves a program from the former toward the latter.
\end{exercise}

\begin{exercise}
Compare a Bash pipeline with a Go program that performs the same sequence of operations on a set of files. For each, identify which intermediate states are explicit, which errors can be observed by the operator, and which failures could leave the filesystem partly changed. Consider in particular what happens when the third of five steps fails.
\end{exercise}

\chapter{Values, Types, and Records}
\label{ch:types}

\section{Structs, Copies, and Zero Values}

A Go program describes its data with types, and the most important composite type for the purposes of this book is the struct, a fixed collection of named fields. The counter of the previous chapter used a struct with a single field, which may have seemed a needless ceremony, but the habit of wrapping state in a named type pays off as soon as the state grows. The artifact package of \golab{} begins with a struct that describes a file that has been found. The struct records three things only, the path, a classification, and the size in bytes at the moment of the walk. It does not record a modification time, an owner, or any property of the content.

\golisting{internal/artifact/types.go}{The types of the artifact package.}{lst:artifact-types}

Three properties of structs govern almost every design decision in the chapters that follow. First, structs are values: assigning one to another variable, passing it to a function, or returning it from a function copies its fields, so the receiver holds an independent value. Second, every type has a zero value that is available without any explicit construction, and for a struct the zero value is the struct whose every field has its own zero value, which is the empty string for a string, zero for a number, and \gocode{nil} for a pointer, slice, or map. Third, a composite literal such as \gocode{Discovered\{Path: p, Kind: k, Size: n\}} names the fields it sets and leaves the rest at their zero values, so that adding a field to a struct does not break existing literals that name their fields.

The zero value is a convenience in many programs and a hazard in programs that wish to say that a value has not been established. A \gocode{Discovered} with an empty path is not a meaningful description of any file, but the language will construct one whenever asked, and it is the responsibility of the surrounding code to refuse to treat it as a discovery. This tension between what the type system can express and what the design requires is one of the main themes of Part~I, and Section~\ref{sec:two-states} returns to it.

\section{Named Types as Vocabularies}

The type \gocode{Kind} is defined as a string with a name of its own, and four constants give it its intended vocabulary: source, output, summary, and other. Naming the type has two benefits. A function that takes a \gocode{Kind} says in its signature what sort of string it expects, and the compiler refuses to pass a plain \gocode{string} variable where a \gocode{Kind} is wanted without an explicit conversion, which catches a class of mistakes in which one kind of text is confused with another.

The benefit has a limit that the book states plainly, because the limit matters to the design. Go has no enumeration type that closes a vocabulary. A \gocode{Kind} is a string underneath, the expression \gocode{Kind("anything")} is a legal conversion, and the compiler will not object to a value outside the four declared constants. The closure of the vocabulary is therefore a convention, maintained by declaring the constants in one place, by routing classification through one function, and by having any consumer that switches on the kind treat an unrecognized value as an error and not as a default. The convention is cheap and effective in a small program, and it is worth being clear that it is a convention. A guarantee would require either a type whose representation cannot be forged from outside its package, which is the technique used for \gocode{Verified} below, or a runtime check at each boundary.

The vocabulary contains a deliberate distinction between a source and a summary. A file ending in \gocode{.summary.md} is classified as a summary, and a file ending in \gocode{.tex} as a source. The classification is made from the name alone and is therefore weak evidence, as Chapter~\ref{ch:traversal} discusses, but the existence of the distinction in the type is what allows later operations to decline to treat a generated summary as if it were the document it summarizes. A summary is an interpretation of evidence and not itself evidence, and a system that can say so in its vocabulary can enforce it in its policies.

\section{Methods and the Package Boundary}

A method is a function with a receiver, written between the keyword \gocode{func} and the method name. The methods on \gocode{Verified} in Listing~\ref{lst:artifact-types} are accessors, and they exist because of the way Go controls visibility. A name that begins with an upper-case letter is exported and may be used by other packages, while a name that begins with a lower-case letter is private to its own package. The fields of \gocode{Verified} are named \gocode{found} and \gocode{sum}, in lower case, so no code outside the artifact package can set them or read them directly. Code elsewhere can obtain a \gocode{Verified} only from the function \gocode{Verify}, which lives in the same package, or as the zero value, and can inspect it only through the accessor methods.

This is a genuine guarantee, and it differs in kind from the conventions discussed above. The Go compiler enforces the visibility rule, so a package that imports \gocode{artifact} cannot write \gocode{artifact.Verified\{sum: "abc"\}} and have it compile. The guarantee is limited in two ways that should be understood exactly. It does not prevent the zero value, which any code may obtain by declaring \gocode{var v artifact.Verified}, and which is why the type carries an \gocode{IsZero} method and why the documentation says that a zero value carries no evidence. The method is specific to \gocode{Verified}, the one type in the module whose zero value would be mistaken for a successful act. The other types do not need it, because their zero values are either harmless, as Chapter~\ref{ch:inscription} shows for \gocode{Valid}, or are already rejected by an existing check. It also says nothing about whether \gocode{Verify} itself is correct, since the type records only that a digest was produced by that function and cannot check that the digest corresponds to the bytes on disk at the moment of use. A file can change after it was verified. What the type can claim is that someone, through the one sanctioned route, read the bytes once and recorded what they hashed to.

\section{Packages, Modules, and the Internal Directory}

The \golab{} module is organized so that each concern lives in a package of its own. The command-line entry point is in \gocode{cmd/golab}, and the reusable machinery lives under \gocode{internal}. The directory name \gocode{internal} has a meaning that the \gocode{go} tool itself enforces: a package under \gocode{internal} can be imported only by code rooted at the parent of that directory. In this module, therefore, \gocode{golab/internal/artifact} can be imported by the command and by other packages of \golab{} but not by any outside module, and the author of the tool is free to change those packages without breaking users who never had access to them. As with the visibility of fields, this is a mechanical guarantee supplied by the toolchain and not a convention.

Dividing the code this way supports the separation of roles introduced in Chapter~\ref{ch:transition}. The package that finds and hashes files knows nothing of the command line, and the package that writes the ledger knows nothing of files. A package that does not import another cannot depend on it, and the Go compiler rejects import cycles, so the dependency structure of the program is a directed acyclic graph that the toolchain checks on every build. The graph for Part~I is simple: the command imports three packages, and none of the three imports another.

\section{Two Types for Two States of Knowledge}
\label{sec:two-states}

The pair \gocode{Discovered} and \gocode{Verified} encodes a distinction that the shell cannot state. A path returned by \gocode{find} is a claim about the existence of a directory entry at one moment and carries no commitment about what the entry contains. A digest of the bytes of a file, computed by reading them, is a claim of a different strength: it supports the later assertion that a file with that digest has exactly those contents, to within the collision resistance of the hash function. By giving the two states different types, the design makes it possible for a function to declare which it needs. An operation that merely lists files can take a \gocode{Discovered}, and an operation that is going to treat the content of a file as established, by recording it in a ledger or comparing it with an earlier digest, can require a \gocode{Verified}, so that a call that skips the reading step does not compile.

The design should not be oversold. It does not make the program correct, and it does not stop a careless author from verifying a file and then ignoring the digest. What it does is move one class of mistake, using an unverified file where a verified one is required, from the category of runtime surprises to the category of compile-time errors, and it records in the program text a commitment about what each stage has and has not established. The same technique recurs throughout the book, and in Part~II it is applied to records and to the stages of the ledger pipeline.

\section*{Exercises}

\begin{exercise}
Add a method \gocode{Valid} to the type \gocode{Kind} that reports whether a value is one of the four declared constants. Decide where in \golab{} the method should be called, and write a test that fails if a new constant is added to the vocabulary without being added to \gocode{Valid}.
\end{exercise}

\begin{exercise}
The zero value of \gocode{Verified} carries no evidence. Write a function that accepts a \gocode{Verified} and returns an error when given the zero value, and explain whether this check belongs in the function, in its callers, or in the type. Consider the cost of forgetting the check in each placement.
\end{exercise}

\begin{exercise}
Replace the exported fields of \gocode{Discovered} with unexported fields and accessors, and observe which parts of the program must change. Argue for or against this change, taking into account that a \gocode{Discovered} makes no claim about the content of the file.
\end{exercise}

\begin{exercise}
Draw the import graph of the \golab{} module as it stands at the end of Part~I. Explain why the Go compiler's refusal of import cycles constrains the placement of the type that represents a ledger record relative to the code that constructs it.
\end{exercise}

\chapter{Traversing, Reading, and Hashing}
\label{ch:traversal}

\section{Walking a Directory Tree}

The operations of the artifact package are in a file of their own, shown below. The first of them, \gocode{Discover}, is the typed descendant of the \gocode{find} command from Chapter~\ref{ch:transition}.

\golisting{internal/artifact/scan.go}{Discovering and verifying files.}{lst:scan}

The function \gocode{filepath.WalkDir} visits every entry beneath a root, calling the supplied function once for each. The function can steer the traversal through its return value. Returning \gocode{nil} continues, returning the special value \gocode{filepath.SkipDir} when positioned on a directory prevents the traversal from descending into it, and returning any other error stops the whole walk and delivers that error to the caller. \gocode{Discover} uses all three behaviours. It skips directories whose names begin with a dot, so that version-control metadata and similar private directories are not mistaken for research material, but it exempts the root itself from the rule so that a caller who names a directory beginning with a dot is not silently handed an empty result. It propagates any error that the walk reports, so that an unreadable subdirectory fails the discovery instead of shrinking it without notice. And it ignores entries that are not regular files, such as symbolic links, sockets, and devices, which is a policy choice stated in the code rather than left to accident.

The choice to fail on an unreadable subdirectory is a decision about policy, and it is worth seeing that it is only one of several defensible decisions. A tool that is used interactively to explore a large tree might prefer to skip what it cannot read and report the omissions at the end. A tool whose purpose is to establish a complete inventory of evidence cannot accept a silently partial inventory, since a later conclusion that something is absent from the directory would then rest on a traversal that never looked. \golab{} takes the second view, and the first exercise at the end of the chapter asks for the alternative.

Two details of the implementation repay attention. The documentation of \gocode{WalkDir} states that it visits entries of each directory in lexical order, but the order of whole paths produced by a depth-first walk is not the same as the lexical order of the complete path strings, because a directory name that is a prefix of a sibling's name interleaves differently. The function therefore sorts the final slice by path, so that the result is deterministic and independent of the details of the walk. Deterministic output matters here because the ledger will record the results of operations, and an operation whose output order varies from run to run would produce records that cannot be compared.

\section{Reading Without Loading}

The function \gocode{Verify} reads a file and records its digest. The central line is the call to \gocode{io.Copy}, which copies from the open file into the hash. That the call is possible at all is a consequence of one of the most characteristic features of Go, the structural interface. The value \gocode{h}, returned by \gocode{sha256.New}, implements the interface \gocode{hash.Hash}, which includes the method \gocode{Write}, and therefore it can be used wherever an \gocode{io.Writer} is expected, with no declaration anywhere that says so. A type satisfies an interface in Go merely by having the required methods. The consequence for this program is that the file is streamed through the hash in blocks, and the program never holds the whole file in memory, so that verifying a ten-gigabyte measurement file costs the same memory as verifying a ten-line note.

The deferred call \gocode{defer f.Close()} schedules the closing of the file to occur when the surrounding function returns, whichever way it returns. Without it, every early return would need its own call to \gocode{Close}, and a forgotten one would leak a file descriptor. A deferred call is run at function exit in the reverse order of its registration, and its arguments are evaluated at the moment of the \gocode{defer} statement. The idiom is so common that it is the first thing a reader of Go code looks for after a successful call to \gocode{os.Open}. For a file opened only for reading, ignoring the error that \gocode{Close} can return is conventional and harmless. For a file opened for writing the same shortcut can lose a failure that matters, and the ledger package, in Chapter~\ref{ch:ledger}, checks the error of \gocode{Close} explicitly for that reason.

The size comparison after the copy is a small defence against a particular failure. The \gocode{Discovered} value carries the size that the file had when the walk observed it, and \gocode{Verify} compares it with the number of bytes actually hashed. If the two differ, the file was modified between discovery and verification, and the function reports the discrepancy and returns no digest. This check is genuinely useful and also genuinely limited. It detects a change in length, and it does not detect a change that preserves length. The gap between a check made at one moment and a use made at another is a general feature of programs that work on a shared filesystem and is known in the security literature as a time-of-check to time-of-use race. \golab{} does not close it and does not claim to. The digest it records attests that the bytes read at a particular time hashed to a particular value, and nothing more.

\section{What a Digest Establishes}

The hash function SHA-256 maps a sequence of bytes of any length to a fixed-length digest of 256 bits, and it is designed so that finding two different inputs with the same digest is computationally infeasible. Under that assumption a recorded digest is strong evidence that a file later found to have the same digest has the same bytes, and equally strong evidence that a file with a different digest differs in at least one byte. That is the entire content of the claim. A digest does not say who wrote the file, whether its contents are true, whether the file is the one the author intended, or whether it is the original from which a copy was made. It is evidence about identity of content and about nothing else.

The point is stated carefully because of a temptation that arises in any tool that records digests. Having computed and stored a digest, it is easy to speak as though the tool had verified the document, and the word invites the further inference that the document is therefore reliable. In \golab{} the type \gocode{Verified} is named for a narrow fact, that the bytes were read and hashed, and the narrow name is deliberate. Any stronger claim would need further evidence, such as a signature by a known key or a comparison with a digest published by an independent party, and the ledger in Chapter~\ref{ch:ledger} is designed so that such evidence can be added without changing the meaning of what is already recorded.

The same reasoning explains why \gocode{Classify} inspects only file names. A name ending in \gocode{.summary.md} is a convention adopted by the operator, and a file with that name might, through error or intent, contain the original text of a source. The classification is documented in the code as a convention, and it is useful because it is cheap and usually right. It is not a verification of the role of the file, and a stricter tool would record, alongside the classification, how it was obtained.

\section{Testing Against the Filesystem}

The tests for the artifact package build their own small directories, and doing so introduces a technique worth learning early. The method \gocode{t.TempDir} returns a fresh directory that the testing framework removes when the test finishes, so a test can create whatever files it needs without risk to the working tree and without cleaning up. The helper function \gocode{write} in the test file creates intermediate directories and writes a file, and calls \gocode{t.Helper} so that a failure inside it is reported at the line of the test that called it. The first test builds a directory containing one source, one output, one summary, and a hidden directory, and checks both that the hidden file is skipped and that each of the others receives the intended classification.

\golisting{internal/artifact/scan_test.go}{Tests for discovery and verification.}{lst:scan-test}

The second test checks the digest of the three bytes \texttt{abc} against the published test vector for SHA-256, which begins \gocode{ba7816bf} and is a standard value in the specification of the algorithm. Testing against a known external value is stronger than testing that the function returns the same thing twice, because it checks the function against an independent source of truth. The third test confirms that a missing file produces an error. The remaining test, added in Chapter~\ref{ch:types}, checks that the zero value of \gocode{Verified} claims nothing.

\section*{Exercises}

\begin{exercise}
Modify \gocode{Discover} so that a directory that cannot be read is recorded in a returned list of omissions and the traversal continues. Change the signature of the function accordingly, and write a test that makes a directory unreadable. Discuss which behaviour a tool used to certify the absence of a file ought to have.
\end{exercise}

\begin{exercise}
Symbolic links are currently ignored. Extend the discovery to report them as a separate kind, without following them, and explain what could go wrong if the walk followed links in a tree that contains a cycle.
\end{exercise}

\begin{exercise}
Construct a scenario in which a file is modified between discovery and verification and the modification is not detected by the size check. Describe a stronger check that would catch it, and state what it would cost.
\end{exercise}

\begin{exercise}
Replace SHA-256 with a different hash function from the standard library, changing as little of the program as possible. Explain which parts of the program had to change, and what that reveals about the role of interfaces in the design.
\end{exercise}

\chapter{Errors as Values}
\label{ch:errors}

\section{The Error Interface}

In many languages a failure interrupts the normal flow of control and travels up the call stack until some enclosing handler catches it. Go makes a different bet. A failure is an ordinary value of the built-in interface type \gocode{error}, whose entire definition is a single method that returns a string, and a function that can fail returns such a value alongside its result. The caller receives the value at the point of the call and must decide, in plain sight, what to do with it. The cost of the arrangement is verbosity, since the test \gocode{if err != nil} appears with a frequency that new readers find tiresome. The benefit is that every place where a failure can occur is visible in the text of the program, and every place where one is ignored is visible as well.

The previous chapters used two constructors for errors. The function \gocode{errors.New} produces an error with a fixed message, and the function \gocode{fmt.Errorf} produces one whose message is formatted from its arguments. A third feature of \gocode{fmt.Errorf} is the one that matters most for the design of \golab{}: the verb \gocode{\%w} wraps another error, so that the new error carries both the added context and the original. In the artifact package, a failure to open a file is reported as an error whose text is of the form \gocode{verify path: open path: no such file or directory}. The outer message says what \golab{} was attempting and the inner one says what the operating system refused, and both are present in the one value.

\section{Sentinel Errors and Inspection}

Wrapping would be of little use if the only way to examine a wrapped error were to search its text. The standard library provides two functions for inspecting the chain. The function \gocode{errors.Is} asks whether any error in the chain is equal to a given target value, and it is the usual way to test for a sentinel, which is a package-level error variable that names a condition. The transition package declares two of them, \gocode{ErrNegative} and \gocode{ErrOverflow}, and the ledger package declares \gocode{ErrBroken}. A caller that wishes to treat an overflow differently from a negative result can write \gocode{errors.Is(err, transition.ErrOverflow)}, and the test succeeds even though the error that \gocode{Add} returns has been wrapped with the numbers involved.

The second function, \gocode{errors.As}, asks whether any error in the chain has a given type and, if so, assigns it to a variable of that type so that its fields can be read. The test \gocode{TestVerifyMissingFileKeepsItsCause}, in Listing~\ref{lst:scan-test}, uses both. It verifies a file that does not exist and then checks that \gocode{errors.Is(err, fs.ErrNotExist)} holds, which shows that the underlying cause survived the wrapping performed by \gocode{Verify}, and that \gocode{errors.As} can recover the \gocode{*fs.PathError} produced by the operating-system layer, whose \gocode{Op} field is \gocode{"open"}. The point of the test is to protect a property that would otherwise erode unnoticed. A later change that formatted an error with \gocode{\%v} in place of \gocode{\%w} would leave every message looking the same to a human reader while severing the chain, and the callers that depend on \gocode{errors.Is} would begin failing to recognize conditions they had previously recognized.

A sentinel is part of the public interface of the package that declares it, in the same way that a function is, and it should be introduced with the same care. Once callers test for \gocode{ErrBroken}, the package has committed to returning it, wrapped or not, in the circumstances that the documentation describes. The sentinels in \golab{} are few and each names a distinct refusal: a transition that would leave the permitted interval on the low side, one that would leave it on the high side, and a ledger whose chain of digests does not hold.

\section{Context and the Cost of Silence}

Error values are most useful when they carry enough context to be acted upon and least useful when they say only that something failed. The functions of \golab{} follow a convention that each layer prefixes the error it returns with the operation it was performing and the object of the operation, so that by the time an error reaches the command it reads as a short narrative. The text \gocode{golab: record 1 digest mismatch: ledger chain broken} shows the pattern. The ledger names the record, the failed check, and the sentinel, and the command adds its own name. A convention of this kind has the property that no layer needs to know the details of another. The artifact package does not know that its errors will be printed by a command, and the command does not know the structure of the errors it prints.

The opposite habit, returning the inner error unmodified, loses the narrative, and the habit of discarding it altogether is worse. The blank identifier \gocode{\_} lets a caller accept a result and ignore the error that accompanied it, and the compiler permits the practice. Static analysers can warn about it, and a project that cares about its failure modes will enable them, but the language places no obstacle in the way. The consequence is that Go provides a convention for handling errors and not a mechanism that forces it. The distinction between reporting a failure and enforcing the boundary the failure was supposed to protect, raised in the fourth exercise of Chapter~\ref{ch:transition}, has its practical form here: a function that returns \gocode{ErrNegative} reports that a transition was refused, and the refusal has an effect on the system only if the caller refrains from using the state it would have produced.

\section{Partial Failure as Data}

Operations that act on many files raise a question that single operations do not, namely what to do when some succeed and some fail. The \gocode{verify} command of \golab{} makes a choice that is worth examining, because it is the choice of a tool that is meant to keep records. It attempts every file, collects the failures instead of stopping at the first, and then appends one record to the ledger that states how many files were found, how many were verified, and how many failed. The record is an aggregate: it holds three counts and does not list the individual files or their outcomes. Per-file results are printed to the terminal and are not stored in the ledger, a limitation that Part~III removes with an optional flag. Only after the record has been written does it return an error if any file failed. The order matters. A command that returned at the first failure would leave the ledger without any record that the attempt had been made, and a command that reported success after any failure would make the record misleading. The chosen order ensures that the ledger always contains an account of what was tried and what happened, and that the exit status of the program tells the operator that the account contains failures.

The behaviour is a small instance of a general principle for systems that must be audited. Failure is information, and a system that records only its successes gives an auditor a distorted picture of what it does. The commands of \golab{} are therefore built so that refusals, failures, and timeouts are written to the ledger in the same form as successes, and the program's exit status is a separate matter, reserved for communicating with the shell.

\section{Errors and Panics}

Go also has a mechanism for conditions that the program cannot reasonably continue from. The built-in function \gocode{panic} stops the ordinary flow of control, runs any deferred calls, and terminates the program with a stack trace unless a deferred function recovers. The runtime raises a panic for errors such as indexing past the end of a slice or dereferencing a nil pointer. The convention of the language is to reserve panics for programmer error and for impossible states, and to use returned errors for anything that can happen in a correctly written program running in an imperfect world, which includes missing files, malformed input, and refused permissions. No code in Part~I calls \gocode{panic}, and the \gocode{main} function of \golab{} does not recover from one. A panic in a tool of this kind is a bug to be fixed, and hiding it would obscure the very evidence that is needed to fix it.

\section*{Exercises}

\begin{exercise}
Change the \gocode{\%w} in the \gocode{Verify} function to \gocode{\%v} and run the test suite. Identify which test fails and explain what property it was protecting.
\end{exercise}

\begin{exercise}
Define a typed error in the transition package that carries the attempted amount and the current counter, and arrange for \gocode{Add} to return it wrapping the appropriate sentinel. Show that a caller can recover both the sentinel and the numbers. Discuss when a sentinel is preferable to a typed error and when it is not.
\end{exercise}

\begin{exercise}
The \gocode{verify} command appends its record to the ledger before returning an error for failed files. Describe a failure of the ledger itself that would defeat this ordering, and propose a policy for the command to follow in that case.
\end{exercise}

\begin{exercise}
Install a static analyser that reports discarded errors, run it against the module, and report what it finds. Then deliberately introduce a discarded error and observe the report. Explain why the analyser is a better description of a policy than of a language feature.
\end{exercise}

\chapter{Running Programs Without a Shell}
\label{ch:processes}

\section{The Argument Vector}

A workflow tool spends much of its life invoking other programs, and the manner of invocation determines how much of the shell's fragility the tool inherits. When a Bash script runs a command, it hands a line of text to an interpreter that splits it into words, expands variables, performs globbing, and interprets quotation, and every one of those steps is an opportunity for a file name containing a space, a quotation mark, or a dollar sign to be read as something other than a name. Go's \gocode{os/exec} package avoids the whole class. The function \gocode{exec.Command} takes the program and each argument as separate strings and passes them to the operating system as the argument vector, with no intervening interpretation. A file called \texttt{my paper (final).tex} is one argument, and it remains one argument however the program is run.

The package \gocode{runner} wraps the facility in a form suited to a tool that must record what happened.

\golisting{internal/runner/runner.go}{Running an external program and classifying the outcome.}{lst:runner}

The function \gocode{Run} starts a program in a given directory, waits for it, and returns a \gocode{Result} that holds the outcome, the exit code, the captured output and error streams, and the elapsed time. Capturing is arranged by assigning \gocode{bytes.Buffer} values to the \gocode{Stdout} and \gocode{Stderr} fields of the command. Because those fields are of type \gocode{io.Writer}, the buffers satisfy them by the same structural rule that allowed a hash to be written to in the previous chapter, and the program is spared the intermediate temporary files that a shell script would use for the same purpose.

\section{Four Endings}

A shell script usually distinguishes two ways in which a command can end, by testing whether its status \gocode{\$?} is zero. A tool that keeps records needs a finer classification, because the operator will ask different questions about different endings. A program that ran and exited with a nonzero status has told us something about its work. A program that was stopped by the tool because it ran too long has told us nothing, and the tool has made a decision on its behalf. A program that could not be started at all, because it does not exist or is not executable, has never run. The type \gocode{Outcome} names the four cases of success, failure, timeout, and failure to start, and the \gocode{switch} in \gocode{Run} sorts every possible ending into exactly one of them.

Go has no sum types, so the four cases are represented as string constants of a named type, with the closure of the set maintained by convention as discussed in Chapter~\ref{ch:types}. The \gocode{switch} makes the order of the tests significant. The timeout is tested before the exit-error case because a process killed by the expiry of the context also yields an error of the exit type, and testing the cause first prevents a killed process from being reported as an ordinary failure with an unusual status.

The function returns a non-nil error only for the fourth ending. This is a design decision with a cost. A caller that tests only \gocode{err} will conclude, wrongly, that a program that ran and failed was a success of the tool, and must also inspect the outcome in the result. The alternative, returning an error for every ending other than success, would have made the signature simpler and the information coarser, and since the ledger needs to record the distinctions, \gocode{Run} reports them in the value and reserves the error for the case in which there is no meaningful value to report. Any such decision should be documented at the function, and the comment above \gocode{Run} states it.

\section{Time and Cancellation}

The timeout is implemented with the \gocode{context} package, which Part~III develops at length. A context is a value that carries a deadline and a signal for cancellation across the boundaries of functions and goroutines. The call \gocode{context.WithTimeout} derives a context that is cancelled when the given duration has elapsed, and \gocode{exec.CommandContext} ties a command to it, so that when the context is done the program is killed. The deferred call to the \gocode{cancel} function releases the resources of the context as soon as \gocode{Run} returns, whether the command finished or not.

The tests for the runner revealed a subtlety that is easy to miss. The timeout test runs the command \texttt{sh -c "sleep 5"} with a limit of a tenth of a second. When the context expires, the kill signal is delivered to the shell, which is the process that was started. The \texttt{sleep} that the shell in turn started is a child of the shell and is not signalled. It still holds the write ends of the pipes through which the output of the shell is being captured, and the standard library's \gocode{Wait} does not return until those pipes are closed. In the first version of the runner, therefore, the test took the full five seconds, because the call was waiting for a process that the tool had already given up on. The remedy is the field \gocode{WaitDelay}, added to \gocode{exec.Cmd} in Go 1.20, which bounds how long \gocode{Wait} will wait for the pipes after the context is done. The runner sets it to half a second, and the test now finishes in a fraction of a second. The episode is a reminder that killing a process is not the same as killing the work it started, and that a tool which supervises programs must decide what happens to descendants. Part~V returns to the question in the setting of containers, where the isolation of process groups and namespaces gives the tool a means of terminating a whole tree.

\golisting{internal/runner/runner_test.go}{Tests for the four endings.}{lst:runner-test}

The tests in Listing~\ref{lst:runner-test} exercise each outcome in turn. They use \texttt{sh} as a convenient stand-in for any program, which ties them to systems that provide a POSIX shell. That dependency is acceptable for a tool developed in a Linux environment, and it is noted here so that it is a decision and not an oversight.

\section{A Shell Pipeline and Its Go Counterpart}

It is instructive to ask what the shell provides that the \gocode{runner} package does not. A pipeline such as \gocode{pdflatex paper.tex | tee build.log} connects the output of one program to the input of another, and the shell arranges the connection. In Go the connection must be made explicitly, either by assigning the \gocode{StdoutPipe} of one command to the \gocode{Stdin} of the next or by passing data through an \gocode{io.Pipe}. The explicitness is a burden for a one-line command and a benefit for a pipeline whose stages must be supervised separately, since each stage then has its own exit status, its own captured error output, and its own timeout, and the failure of one is visible instead of being masked by the status of the last. The shell offers the option \gocode{pipefail} to approximate this, and Go makes it the default shape of the code.

The bounded character of the \gocode{runner} package is deliberate. It runs one program, waits for it, and reports what happened. Chaining, retrying, and running programs in parallel are operations on top of it, and Part~III builds them from goroutines and channels so that the supervision of a single program, once settled, is not reopened.

\section*{Exercises}

\begin{exercise}
Write a test that runs a program with an argument containing spaces, quotation marks, and a dollar sign, and shows that the program receives it unchanged. Then write the corresponding Bash command and describe what quoting is required to achieve the same result.
\end{exercise}

\begin{exercise}
Remove the assignment to \gocode{WaitDelay} and run the timeout test. Measure how long it takes and explain the result in terms of the processes and pipes involved. Then change the test command to \texttt{exec sleep 5} and explain why the behaviour differs.
\end{exercise}

\begin{exercise}
Extend \gocode{Result} so that it also reports the signal that terminated a process, where there was one, and add a fifth outcome for a program killed by a signal other than the tool's own. State what information the operator gains from the distinction.
\end{exercise}

\begin{exercise}
Implement a function that runs two programs as a pipeline, connecting the output of the first to the input of the second, and returns a separate result for each. Decide how the function should behave when the second program exits before the first has finished writing.
\end{exercise}

\chapter{A Ledger of Actions}
\label{ch:ledger}

\section{What a Ledger Is For}

A tool that finds files, hashes them, and runs programs on them produces consequences, and an operator who returns to the tool after a week will want to know what it did. Terminal output is gone by then, and the state of the filesystem records only the net effect of the actions and not the sequence. A log file is the usual remedy, but an ordinary log has a weakness that matters in a setting where the point of keeping records is to rely on them later: any line can be changed or removed, and nothing in the file reveals that the change took place. The ledger of \golab{} is an append-only sequence of records in which each record includes the digest of the record before it, so that a change to an earlier line breaks the chain and can be found by recomputation.

The ledger records events and does not judge them, and it is worth being exact about what that sentence covers. In Part~I the ledger receives one record for each \gocode{verify} and each \gocode{run} command, and nothing for \gocode{discover}, which is read-only. A \gocode{verify} record holds counts and not per-file results. No record names the person or process that issued the command, since the tool has no notion of identity. The ledger also does not accept every write unconditionally: \gocode{Append} refuses to extend a chain that already fails the check, which is a refusal made by the ledger itself and is deliberate. It states that a verification was performed, that a program was run and ended in a certain way, and that these happened in this order, and it takes no position on whether they were permitted. That separation between recording and authorizing is one of the central distinctions of the book, and in Part~I it is observed by construction, since the program has no policy beyond a few classification rules and the admissibility checks of the transition package. Part~II introduces policy as a separate component and then raises the question of what a ledger of events should say about refusals.

\section{The Record and Its Digest}

A record has a sequence number, a timestamp, an action, a subject, a map of details, the digest of its predecessor, and its own digest. The structure is declared with field tags that control its encoding as JSON, and each record occupies one line of a file, a format commonly called JSON Lines, which has the convenient property that appending a record never requires rewriting what is already there.

\golisting{internal/ledger/ledger.go}{The ledger package.}{lst:ledger}

Let the records be $r_1, r_2, \ldots, r_N$ and let $\operatorname{enc}$ denote the JSON encoding of a record with its digest field set to the empty string. The digest of record $n$ is defined by
\[
d_n = H\bigl(\operatorname{enc}(r_n)\bigr), \qquad d_0 = \varepsilon,
\]
where $H$ is SHA-256 and $\varepsilon$ the empty string, and the field of $r_n$ that records its predecessor holds $d_{n-1}$. Because $d_{n-1}$ is part of the encoded content of $r_n$, the value $d_n$ depends on every earlier record, through the chain of predecessors. The function \gocode{digest} in Listing~\ref{lst:ledger} computes the definition directly. It relies on a property of the JSON encoder that the documentation guarantees, namely that the keys of a map are written in sorted order and that the fields of a struct are written in the order of their declaration, so that the same record always encodes to the same bytes. Without that determinism the digest would not be reproducible, and a ledger written on one machine could not be checked on another.

A change to a record is reported at that record's own sequence number, because its recomputed digest no longer agrees with the stored one, and the linkage test at the next record is reached only if the digest was rewritten. The function \gocode{Check} verifies three things for each record in turn: that its sequence number equals its position, that its predecessor field equals the digest of the preceding record, and that its own digest equals the recomputed value. It reports the first failure, wrapped around the sentinel \gocode{ErrBroken}, and the wrapping carries the sequence number so that an operator is told where to look. The function \gocode{Append} loads the existing ledger, checks it, and refuses to extend one that fails the check, because appending to a broken chain would give the damage a respectable appearance. The new record is written with the file opened in append mode, and the error of the final \gocode{Close} is examined, since a failure to flush a write is reported there and not at the call to \gocode{Write}. The listing is the final form of the package, which Part~VI revised in two ways that this chapter does not discuss: a write that fails part-way is cut back so that no fragment of a record remains, and the record is flushed to stable storage before the function returns. Chapter~\ref{ch:faults} gives the reasons.

\section{What the Chain Establishes}

The guarantee that the chain gives is narrower than it can appear, and it is worth stating exactly.

\begin{proposition}\label{prop:chain}
Suppose a ledger of $N$ records satisfies \gocode{Check}, and suppose $H$ is collision resistant. If a single record $r_k$ with $k < N$ is replaced by a different record while all others are unchanged, then \gocode{Check} fails. If a trusted party holds the head digest $d_N$, then any ledger that satisfies \gocode{Check} and has head digest $d_N$ has the same records as the original, in the same order.
\end{proposition}

\begin{proof}
Let $r_k'$ be the replacement for $r_k$ with $k<N$. If $r_k'$ keeps the old digest field, then the digest recomputed from its content differs from the stored one unless the two records encode identically, which contradicts their being different. If $r_k'$ stores a digest recomputed from its own content, that digest differs from $d_k$ by collision resistance, and record $r_{k+1}$ still stores $d_k$ as its predecessor, so the linkage test fails at $k+1$. For the second claim, a ledger that satisfies \gocode{Check} is determined by its records, and each digest is a function of everything before it, so two such ledgers with the same $d_N$ must agree record by record, or else a collision exists.
\end{proof}

The proposition has two limits that follow from its own statement. The first is that the check alone does not protect the final record or the length of the ledger. A person who can rewrite the file can alter the last record and recompute its digest, or can delete trailing records, and the shortened file will still satisfy \gocode{Check}. The second is that the check establishes internal consistency and nothing about authenticity: an adversary who controls the whole file can produce a different ledger with perfectly consistent digests from beginning to end. The protection against both is the one the proposition names, which is a copy of the head digest held somewhere the adversary cannot reach, such as a printed note, a message sent to another party, or a record in a different system. The ledger is therefore best described as tamper-evident relative to an external anchor. It is not tamper-proof, and the version in Part~I does not provide the anchor. The tool does not yet have a command that prints the head digest, and the first exercise adds one. Where the digest is stored is the operator's decision, and the examples above are illustrations and not features of \golab{}.

Two further limits belong to the implementation. Timestamps come from the clock of the machine, are truncated to whole seconds, and are not authenticated, so the ledger shows the order in which events were recorded and not trustworthy times. The subject of a run record is the program and its arguments joined with spaces, which loses the original quoting, as the transcript below shows, and is therefore a description for a human reader and not a command that can be re-executed.

\golisting{internal/ledger/ledger_test.go}{Tests for the ledger chain.}{lst:ledger-test}

The tests of Listing~\ref{lst:ledger-test} check the three behaviours that matter. The first appends three records and confirms that the chain holds when reloaded. The second corrupts the file by changing one character of the first record and confirms that \gocode{Check} reports \gocode{ErrBroken} and that \gocode{Append} refuses to extend the damaged ledger. The third confirms that a missing file is read as an empty ledger, which is the state in which the first append creates the file.

\section{The Command and a Recorded Session}

The command-line program connects the pieces. Its \gocode{main} function does nothing except pass the arguments to a function named \gocode{run} and exit with the status that \gocode{run} returns, so that the decision about the process exit status is made in one place and the rest of the program is testable. The complete source is in Appendix~\ref{app:main}. Each subcommand follows the same shape, in which it performs an operation, appends a record of the outcome to the ledger, and only then returns an error if the operation failed, which is the ordering argued for in Chapter~\ref{ch:errors}.

The following transcript is the output of an actual session. The directory \texttt{papers} contains a text file, a \TeX{} file, and a summary of the second. The session discovers the files, verifies them, runs two programs of which one fails, audits the ledger, and then alters the first line of the ledger and audits again.

\lstinputlisting[style=shellstyle,caption={A recorded session of \golab{}.},label={lst:session}]{session.txt}

Several features of the transcript carry the argument of the chapter. The kinds in the first listing show the classification at work, with the summary distinguished from the sources. The digest of \texttt{a.txt} is the standard value for the three bytes \texttt{abc}. The failed program is recorded in the ledger with its outcome and exit status, and the command exits with status 1 so that a surrounding shell script can detect the failure. The audit reports that the chain holds until the first line is altered by changing a single letter, after which the audit fails and, more significantly, the next attempt to run a program is refused until the ledger is repaired. A tool that continued to extend a ledger it knew to be damaged would be undermining the one property that gives the ledger its use.

\section{Where Part I Ends}

The tool at this point can find files, state what it has read of them, run programs on them, and keep a chained record of what it did, and each of these capacities has a type or a package that names it. The four responsibilities of Chapter~\ref{ch:transition} have homes: the command in \gocode{cmd/golab}, the operations in \gocode{artifact} and \gocode{runner}, the record in \gocode{ledger}, and a rudimentary policy in the classification rules and the transition checks. What the tool lacks is also clear. It has no explicit policy that decides whether an operation is permitted, and the ledger cannot say that something was refused for reasons of policy. It cannot safely run two operations at once, since two processes that append at the same moment can each read the same predecessor and write records with the same sequence number, a fault that \gocode{Check} would detect afterward and that nothing prevents. It does not anchor the head of its ledger anywhere outside the file. These absences are the agenda of the later parts of the book, and the plan for them follows this chapter.

\section*{Exercises}

\begin{exercise}
Add a command \gocode{head} that prints the sequence number and digest of the last record, and write a test. Then describe a procedure by which an operator could use the output to detect the deletion of trailing records, and state what the procedure assumes about where the output is kept.
\end{exercise}

\begin{exercise}
Run two \gocode{golab run} commands at the same instant in a loop until the ledger is damaged, and report how the damage appears to \gocode{Check}. Propose a locking strategy based on a lock file, and identify a failure of that strategy.
\end{exercise}

\begin{exercise}
Extend the record so that a verification event includes the digests of the files it verified, and consider the size of the resulting record for a directory of ten thousand files. Discuss whether the digests belong in the ledger or in a separate manifest whose own digest is recorded in the ledger.
\end{exercise}

\begin{exercise}
Design, without implementing it, a record for a refusal, that is, for a command that was rejected before any operation was attempted. State which fields it needs, what the subject should be, and why the ledger of Part~I has no honest way to produce it.
\end{exercise}


\part{Governed Workflows}
\chapter{Interfaces and Composition}
\label{ch:interfaces}

\section{Why Policy Is a Separate Thing}

At the end of Part~I the tool could do things and record that it had done them, but it had no way of saying that a thing should not be done. The only rules in the program were scattered through the code, as the classification of file names, the bounds of the counter, and the refusal to extend a broken ledger. A rule that lives inside the function it governs cannot be varied without editing the function, cannot be shown to an auditor without reading the function, and cannot be combined with other rules except by writing a larger function. The design question of Part~II is how to move rules out of the operations and into values that can be named, passed around, tested on their own, and combined.

Go answers the question with the interface, and the answer is characteristic of the language. An interface in Go is a set of method signatures, and a type satisfies it by having those methods, with no declaration of intent anywhere. The standard library uses this heavily. The hash of Chapter~\ref{ch:traversal} satisfied \gocode{io.Writer} without knowing the interface existed, and the buffers of Chapter~\ref{ch:processes} did the same. The consequence for design is that the consumer of a behaviour defines the interface it needs, in its own package, and any type anywhere that happens to have the right methods can be supplied. The package of this chapter defines the single interface through which every rule of \golab{} is expressed. The interface has one method, named \gocode{Decide}, which takes a record and returns a verdict.

\golisting{internal/pipeline/policy.go}{Verdicts, the policy interface, and combinators.}{lst:policy}

\section{Verdicts as an Ordered Vocabulary}

A policy examines a record and returns a verdict, and the verdict is an integer type whose five constants are declared with the keyword \gocode{iota}. Within a constant declaration, \gocode{iota} starts at zero and increases by one on each line, so \gocode{Allow} is $0$, \gocode{Unscoped} is $1$, and so on up to \gocode{Refuse}, which is $4$. The choice of an integer representation, in place of the string constants used for \gocode{Kind} and for the outcomes of the runner, is deliberate, because verdicts are ordered by severity and the order is used in the combinator below. The order is a fact about the vocabulary and is stated once, in the declaration, by the placement of the constants.

The method \gocode{String} on \gocode{Verdict} makes the values printable by name. The \gocode{fmt} package checks whether a value has a method with that signature, which is the interface \gocode{fmt.Stringer}, and uses it when formatting, so that a verdict in an error message appears as \texttt{refuse} and not as the number four. Again no declaration connects the type to the interface. The unexported method \gocode{disposition} maps each verdict to the disposition that an admitted record will receive, and it ends in a default branch that treats any unrecognized value as a refusal. A system that decides what is permitted should fail closed, meaning that when it meets something outside its vocabulary it denies and does not allow, and the branch is where that principle is applied to a conversion.

The reader will notice that the policy returns a \gocode{Verdict} and not a \gocode{Disposition}, although the two are closely related and the second has seven values to the first's five. The difference is the point of the next chapter. Two of the seven dispositions, duplicate and malformed, are facts about the record that the pipeline determines for itself, and a policy has no business asserting them. By giving the policy a smaller vocabulary, the type system prevents a policy from claiming that a record is a duplicate, and the compiler enforces the restriction because a function that returns \gocode{Verdict} has no way to produce the other two.

\section{Functions with Methods}

Writing a new named type for every rule would be tedious, and Go provides a compact alternative that surprises readers coming from other languages. Any named type can have methods, including a named function type. The declaration of \gocode{PolicyFunc} says that a \gocode{PolicyFunc} is a function from a record to a verdict, and the method \gocode{Decide} on it simply calls the function. The effect is an adapter: any function with the right signature can be converted to a \gocode{PolicyFunc} and thereby becomes a \gocode{Policy}. The standard library's \gocode{http.HandlerFunc} is the best-known instance of the idiom, and Part~V will use it.

The constructors \gocode{DenyAuthors}, \gocode{RequireScope}, and \gocode{GateContaining} each return a policy built from a closure, a function that refers to variables of the function that created it. The closure in \gocode{DenyAuthors} captures a map of denied names, built once at construction and consulted on every call, so that the cost of building the set is paid when the policy is made and not each time a record is judged. The variables captured by a closure live as long as the closure does, and the garbage collector of Go makes that safe without the programmer managing it. One caution applies to every such construction. A closure that captures a map or slice shares it with whoever else holds it, so a policy that captured a map handed in by the caller would change behaviour if the caller later modified the map. The constructor here copies the names into a map of its own, and the habit of copying at the boundary appears again in the next chapter.

\section{Composition and What It Preserves}

The function \gocode{AllOf} combines any number of policies into one, and its rule is that the most restrictive verdict wins. Because the verdicts are ordered integers, the rule is the arithmetic maximum, and the implementation is a loop that keeps the largest value seen. An empty combination allows everything, which is a decision with consequences. A caller who builds a combination from a list that happens to be empty gets a policy that admits every record, and a system intended to deny by default would need to place an explicit refusal at the bottom of the list. The behaviour is documented at the function, and the command-line tool uses it deliberately, since running \gocode{golab admit} with no flags is a request to apply no rules.

The following proposition records what the combinator guarantees, and the guarantees are the reason the design can be relied on when policies are assembled from several sources.

\begin{proposition}\label{prop:allof}
Let $p_1,\ldots,p_n$ be policies and let $p = \texttt{AllOf}(p_1,\ldots,p_n)$. Then for every record $r$, $p(r) = \max\{\mathrm{Allow},\, \hat p_1(r), \ldots, \hat p_n(r)\}$, where $\hat p_i(r)$ is $p_i(r)$ if that is a declared verdict and $\mathrm{Refuse}$ otherwise. Consequently (a) the result does not depend on the order in which the policies are listed or on any repetition among them, (b) for any further policy $q$, $\texttt{AllOf}(p_1,\ldots,p_n,q)(r) \ge p(r)$, and (c) a record is admitted under $p$ only if it is admitted under each $p_i$ taken alone.
\end{proposition}

\begin{proof}
The formula is the definition of the loop. Part (a) holds because the maximum of a finite set is unchanged by reordering or repeating its members. Part (b) holds because adding a member to a set cannot lower its maximum. For (c), a record is admitted exactly when its verdict is \gocode{Allow}, the lowest value, so $p(r)=\mathrm{Allow}$ forces every $\hat p_i(r)$ to be \gocode{Allow}, and a declared $p_i(r)$ equal to \gocode{Allow} is then the verdict of $p_i$ alone.
\end{proof}

The proposition is a statement about a join-semilattice in which the policies are combined by taking the maximum, and its practical content is that adding a rule to a combination can only make the system stricter. A reader who wants to know why a record was admitted need only check that every rule individually permits it. The property is checked by tests in the next chapter, which shuffle the order of the policies and confirm that the dispositions do not change, and which add policies one at a time and confirm that the number of admitted records never rises.

The first draft of \gocode{AllOf} did not contain the clause about undeclared verdicts. Its loop began at \gocode{Allow} and kept any larger value, which silently discarded a verdict below \gocode{Allow}. A policy that returned $-1$ was refused when used alone, because the conversion to a disposition falls through to the refusing default, and was ignored when used inside \gocode{AllOf}, because the maximum began at zero. A verdict above \gocode{Refuse}, such as $99$, never had this problem: it won the maximum and was then refused by the same default. The flaw was confined to values below \gocode{Allow}. The inconsistency was found while writing the proof above, when the formula for $p(r)$ could not be made to hold for all integer values, and the repair is the call to \gocode{valid} that sends any undeclared verdict to \gocode{Refuse}. The case shows a general hazard of integer vocabularies in a language without closed enumerations, and it is the reason the tests include verdicts outside the declared range on both sides.

\section*{Exercises}

\begin{exercise}
Write a policy constructor \gocode{AnyOf} that returns the \emph{least} restrictive verdict of its arguments. State the algebraic properties that correspond to those of Proposition~\ref{prop:allof}, and explain why a system that combined independent authorities with \gocode{AnyOf} would be harder to audit than one that used \gocode{AllOf}.
\end{exercise}

\begin{exercise}
Define a policy as a struct with a method instead of a closure, for instance a \gocode{Denylist} type holding a map, and show that it satisfies \gocode{Policy}. Compare the two styles with respect to inspection: which makes it easier for a program to print the rules it is applying?
\end{exercise}

\begin{exercise}
Make \gocode{AllOf} with no arguments refuse everything, and identify every place in the module whose behaviour changes. Argue for or against the change.
\end{exercise}

\begin{exercise}
Demonstrate the aliasing hazard described in this chapter. Write a constructor that captures a caller's slice without copying it, show a test in which later modification of the slice changes the behaviour of an already constructed policy, and then repair the constructor.
\end{exercise}

\chapter{Inscription and Validation}
\label{ch:inscription}

\section{Three Questions That Are Easy to Merge}

Chapter~\ref{ch:transition} observed that a single function can answer, without saying so, three different questions about a proposed change: whether it is well formed, whether the actor is entitled to request it, and whether it is made durable. The distinction is easy to lose because a small program answers all three in one place and nothing breaks. It matters as soon as a system must explain itself, because the three questions have different answers to the question of who is responsible. A record can be perfectly well formed and still not be permitted. A record can be permitted and not yet stored. A record can be stored, receipted, and preserved forever, and still never be treated as true.

The pipeline of this part separates the questions into stages and gives each stage a type. A record is first \emph{inscribed}, which means stored and receipted, and which involves no judgment of any kind. The set of inscriptions is then \emph{validated}, which checks the structure of the whole set without reference to anyone's rules. The valid set is then \emph{admitted} under a policy, which assigns each record a disposition. The admitted records are \emph{projected} onto a surface, meaning that the records belonging to a given scope are selected, and the projection is \emph{activated}, which makes it live provided that it is internally coherent. The five stages follow the staging described in the author's specification of a public claim graph. The code of this chapter and the next is a teaching model of that staging and does not implement the specification: it has no signatures, no argument semantics, and no claim to the formal properties that the specification proves. Where the text speaks of a property, it means the property of this small program, checked by this program's tests.

\golisting{internal/pipeline/record.go}{Records, receipts, and inscription.}{lst:record}

\section{Inscription}

The pipeline of Part~II is a separate chain of types from the artifact types of Part~I. \gocode{Inscribe} takes a \gocode{Record}, which is a claim with an author and a body, and not a file; it returns an \gocode{Inscribed}, which is not a \gocode{Verified}; and it does not write to the ledger. The two chains meet only in the command, which records the outcome of an \gocode{admit} run as a ledger event. A \gocode{Record} carries an identifier, an author, a scope, a body, and a list of identifiers of records it refers to. The function \gocode{Inscribe} encodes the record as JSON, hashes the encoding with SHA-256, and returns an \gocode{Inscribed} value holding the record and the resulting receipt. The signature is the whole of its specification: it takes a record and returns an inscription, it has no error result, and it has no policy parameter. A record with no author, no body, and no identifier is inscribed as readily as a perfect one. This is an intended property and not an oversight, and it expresses the principle that storing is not endorsing. An archive that refused to store a document it disapproved of would no longer be an archive, and a ledger that silently dropped malformed records would conceal exactly the events an auditor most wants to see.

The receipt is deterministic, since the encoding of a struct follows the order of its fields and the encoding of a string slice follows the order of its elements, and so the same record always yields the same receipt, on any machine. The function \gocode{copyRecord} makes a fresh copy of the \gocode{Refs} slice before the record is stored. The copy answers a hazard that is specific to Go's slices. A slice is a small value that refers to an underlying array, so assigning a slice, or placing it in a struct, shares the array and does not duplicate it. Had \gocode{Inscribe} stored the record as given, a caller who later modified the slice it passed in would change the contents of an inscription after its receipt had been issued, and the receipt would no longer describe the record it was supposed to attest. The test \gocode{TestInscribeCopiesRefs} fixes the behaviour. The accessor \gocode{Record} copies on the way out for the same reason.

\section{Validation}

Validation examines a set of inscriptions as a whole and asks whether it is structurally sound. A sound set has three properties: every record has a non-empty identifier, no identifier is attached to two different contents, and the references among the records form a graph with no dangling references and no cycles. The listing below shows the validation function together with the later stages, which the next chapter takes up.

\golisting{internal/pipeline/stages.go}{Validation, admission, projection, and activation.}{lst:stages}

The function \gocode{Validate} builds a map from identifiers to inscriptions and checks the first two properties as it goes. Inscribing the same record twice is treated as idempotent, since the receipts agree, and so a set can be assembled from overlapping sources without being rejected for containing a repeat. An identifier is therefore allowed to appear more than once, provided that every appearance has identical content. Two different records under one identifier are an inconsistency. The check for dangling references runs next and the check for cycles last, using a depth-first search in which each identifier is white while unvisited, grey while its references are being explored, and black once all of them are done. Meeting a grey identifier during the search means that a reference leads back to a record whose own references are still being followed, which is a cycle. The three-colour scheme is the standard method, and it runs in time proportional to the number of records and references.

A detail of the implementation deserves comment because it affects reproducibility. The iteration order of a Go map is deliberately randomized by the runtime, so a loop over the map would visit the records in a different order on each run, and a validation that reported the first error it found would report a different one from run to run. The function therefore collects the identifiers, sorts them, and iterates over the sorted slice, so the error returned for an invalid set is the same every time and the valid set is stored in identifier order. A reader who has written shell scripts that loop over a glob will recognize the concern, since there too the output order is a property that must be established and cannot be assumed.

\section{Validity Does Not Mention Policy}

The central claim of this chapter is that validation, and inscription before it, are independent of any policy. The claim is expressed in the signatures: \gocode{Inscribe} takes a record and \gocode{Validate} takes a slice of inscriptions, and neither takes a \gocode{Policy}. The compiler therefore guarantees that no policy can be passed to them, and a function that cannot receive a policy cannot vary its behaviour according to one.

\begin{proposition}\label{prop:independence}
For every finite set of records $R$ and every pair of policies $p$ and $q$, the receipts produced by inscribing $R$ and the result of validating the inscriptions are the same whether the set is subsequently admitted under $p$ or under $q$.
\end{proposition}

\begin{proof}
Inscription and validation are functions of the records alone. Admission is applied to the output of validation and is not an input to it, and no stage feeds a policy's result back to an earlier stage. Hence the receipts and the validity of $R$ are determined before either policy is consulted.
\end{proof}

The proof is short because the design makes it short, and the honesty of the claim depends on a qualification that the language forces. Go does not enforce purity. A function that takes no policy parameter could still read a package-level variable that holds a policy, or the clock, or a file, and the signature would not reveal it. The compiler's guarantee is the narrow one that no policy is passed in, and the broader statement that nothing else influences the result rests on the code of the two functions, which can be read in a minute, and on the test \gocode{TestPolicyDoesNotAffectInscriptionOrValidity}. The test builds the same records, runs them through admission under three quite different policies, and checks that the receipts and the size of the valid set are identical each time. It checks instances of the proposition and does not prove it. The distinction between an argument that covers all cases and a test that covers some is one the book keeps in view throughout.

\section{The Zero Value, Again}

The type \gocode{Valid} follows the pattern of \gocode{Verified} in having unexported fields, so that only \gocode{Validate} can build a non-empty one. Its zero value, however, differs in an instructive way. The zero \gocode{Verified} claimed to have read a file and had not, so it was dangerous and needed an \gocode{IsZero} guard. The zero \gocode{Valid} is an empty set of records, and an empty set satisfies every structural property vacuously, so it is a legitimate valid set and passing it on is harmless. A type whose zero value is meaningful and safe needs no guard, and designing types so that the zero value is a sensible state is an old Go proverb that this contrast makes concrete.

\section*{Exercises}

\begin{exercise}
Change \gocode{Validate} so that a reference to an unknown record is allowed and treated as a pending reference. Identify which later stage must then change, and what the proposition on policy independence still guarantees.
\end{exercise}

\begin{exercise}
Remove the sort from \gocode{Validate} and run the tests repeatedly with the \gocode{-count} flag. Report which tests become unreliable and explain why in terms of map iteration.
\end{exercise}

\begin{exercise}
The receipt covers the record but not the time at which it was inscribed. Add an inscription time without breaking the determinism of receipts, and decide whether the time should be part of the receipt.
\end{exercise}

\begin{exercise}
Write a function that, given a valid set, returns the records in an order such that every record follows the records it refers to. Explain why such an order exists exactly because the set is acyclic.
\end{exercise}

\chapter{Admission, Projection, and Refusal}
\label{ch:admission}

\section{Admission}

Admission is the first stage at which a policy is consulted, and it is where a valid set of records receives its dispositions. The function \gocode{Admit}, in Listing~\ref{lst:stages}, visits the records in identifier order and assigns each one of seven dispositions. Two of them are determined by the pipeline itself. A record with no author or no body is malformed, and a record whose author, scope, and body repeat those of an earlier record in identifier order is a duplicate. Only a record that is neither is shown to the policy, and the policy's verdict is converted to one of the remaining five dispositions: admitted, unscoped, gated, quarantined, or refused.

The ordering of these tests encodes a decision. Malformed records are never shown to a policy, so a policy is never asked to reason about a record that lacks the fields it would inspect, and a policy written without defensive checks cannot be made to fail by a record with an empty body. A duplicate is likewise detected before the policy runs, so the policy is not asked to judge the same content twice. The consequence to be aware of is that duplicate detection depends on the order of identifiers, since the earlier record in that order is the original and the later is the repeat. Two records with identical content but different identifiers are therefore treated asymmetrically, and an author who wishes a particular one to count as the original must choose identifiers accordingly. The behaviour is deterministic and documented, and it is a convention of this model and not a law.

A record that the policy refuses still registers as seen, so that a later identical record is reported as a duplicate of a refused record. This is a small instance of a general question for any system that keeps state about what it has judged: whether the memory of a judgment should include the judgments that went against the subject. The choice here is the simpler one, and the first exercise asks for the alternative.

\section{Seven Dispositions, Five Verdicts}

The split between a seven-valued disposition and a five-valued verdict deserves a second look, because it is an example of using types to divide authority. The pipeline holds authority over what is structurally wrong or redundant, and the policy holds authority over what is permitted, and each is confined to its own vocabulary. A policy author cannot declare a record malformed in order to hide it, since the type of the policy's result has no such value, and the pipeline's author cannot make a record refused for policy reasons without a policy being consulted. The separation does not decide who deserves to hold either authority. It makes the boundary visible in the types, so that anyone who wishes to move it has to change a declaration that is seen in review and cannot do so by an incidental edit to a function body.

\section{Projection and Activation}

A projection selects the admitted records visible on one surface, where a surface is identified by a scope prefix. The function \gocode{Project} takes only records whose disposition is admitted and whose scope begins with the given prefix, and it copies them so that the projection owns its data. Activation then makes the projection live, subject to one check that cannot be done earlier. A projected record may refer to a record that the projection left out, either because that record belongs to another scope or because it was not admitted, and an active record whose foundations are invisible on its own surface is an orphan. The function \gocode{Activate} returns \gocode{ErrOrphaned} if any projected record has such a reference. The check is a single pass over the direct references of the projected records and is not recursive, because if every direct reference of every projected record lies inside the projection, then so does every indirect one.

The need for the check shows why the stages are distinct. Validation found every reference intact, since the full set was closed and acyclic. Admission may then have refused or gated some of the referred-to records, and projection may have separated the rest by scope, and the coherence of the set was destroyed by operations that were each individually reasonable. No earlier stage could have detected this, because the damage did not exist until later stages acted. A system that validated once at the start and never again would present orphans as though they were sound.

\golisting{internal/pipeline/pipeline_test.go}{Tests for the pipeline.}{lst:pipeline-test}

\section{What the Tests Check}

The tests in Listing~\ref{lst:pipeline-test} are organized around the claims of the last three chapters. The table in \gocode{TestValidate} covers the success case and each failure the validator can report, including the repeat that is accepted as idempotent and both the self-cycle and the two-record cycle. The test \gocode{TestAdmitDispositions} builds a set of six records, chosen so that five of the seven dispositions are exercised (the unscoped disposition appears in the session below, and the quarantine only in the combinator tests), and compares the result with the expected dispositions in full. The record refused for its author and also outside the required scope is expected to be refused, which checks that the more severe verdict prevails. The tests \gocode{TestAllOfIsOrderIndependent} and \gocode{TestAddingAPolicyNeverAdmitsMore} check parts (a) and (b) of Proposition~\ref{prop:allof} on instances, and \gocode{TestUnknownVerdictFailsClosed} checks both the direct and the combined path for verdicts outside the declared range. Finally \gocode{TestActivateDetectsOrphans} builds the situation described above, in which one record in one scope refers to a record in another, and confirms that activation fails for the narrow scope and succeeds for the wide one.

Tests of this kind are evidence of a specific character. Each one confirms that for the particular records and policies it constructs the program behaves as the proposition says it must. None of them shows that the proposition holds for all records and all policies, and the reader should not read the passing suite as a proof. The propositions of the chapters are proved in the text, and the tests guard against the code departing from the text. The first draft of \gocode{AllOf} passed every test that existed when it was written and was nevertheless inconsistent with the proposition, and it was the attempt to prove the proposition and not any test that revealed the gap. The experience is a fair summary of the relation between the two activities.

\section{Recording Refusals}

Part~I ended with a limitation: the ledger could record that an operation was performed and had no honest way to record that one was refused. The \gocode{admit} subcommand closes the gap in the following way. It reads a file of records, inscribes and validates them, builds a policy from its flags, admits the records, prints one line per record with its disposition and the first twelve characters of its receipt, and appends a single ledger record whose details are the number of records and the number that received each disposition. A validation failure is recorded too, under a distinct action name, with the stage and the error, so that an attempt to admit an inconsistent set leaves a trace. If activation fails the failure is placed in the details of the ledger record before the error is returned, in keeping with the ordering argued in Chapter~\ref{ch:errors}.

\lstinputlisting[style=shellstyle,caption={A recorded session of the \gocode{admit} command.},label={lst:session2}]{session2.txt}

In the first command of the transcript the records are judged under three rules at once. The record by the denied author is refused, the one in another scope is unscoped, the repeat is a duplicate, the one without an author is malformed, and the one containing the gated phrase is held for review. Exactly one record is admitted, and it is the only one active in the scope. The second command applies no rules, admits everything that is not malformed or a duplicate, and then projects the second scope, in which the record that remains refers to a record outside it. Activation fails, the command exits with a nonzero status, and the ledger records both runs.

The transcript also shows what the ledger does not yet record. The ledger entry holds counts and not the individual dispositions, so an auditor can see that one record was refused but not which. It holds no description of the policy that was applied, so a run with one set of flags is indistinguishable in the ledger from a run with another that happened to yield the same counts. And the receipts that the command prints are not stored. These omissions are deliberate simplifications for the sake of the chapter and are the subject of the exercises. A ledger that is to support an audit needs to record the policy as well as the outcome, because a disposition is a function of both the record and the rules, and a record of the first without the second cannot be reproduced.

\section{Where Part II Ends}

The tool now distinguishes four things that the shell blurs and that Part~I could only name: a record that has been stored, a record that is structurally sound as part of a set, a record that has been judged under a policy, and a set of records that is coherent on a surface. Each has its own type, each stage has its own function, and the dependence of each on the one before it is visible in the signatures. What the tool lacks is also clear. It runs one stage at a time on one thread, and it does not record the policy it applied. It treats a file of records as the whole world, and it offers no way for several operators to submit records at once. Part~III takes up concurrent execution, beginning with the worker pool that verifies many files at the same time and the single goroutine that will own the ledger.

\section*{Exercises}

\begin{exercise}
Change \gocode{Admit} so that a refused record does not register as seen, and show an input on which the dispositions differ from those of the present version. Discuss which behaviour is preferable for an audit.
\end{exercise}

\begin{exercise}
Extend the ledger record written by \gocode{admit} to include the full list of identifiers and dispositions, and measure the size of the record for ten thousand records. Propose a representation in which the ledger holds only a digest of the full list and the list is kept elsewhere, and state what the operator must do to check the digest.
\end{exercise}

\begin{exercise}
Add a description of the policy to the ledger record, produced by a \gocode{Describe} method that each policy implements. Consider whether \gocode{PolicyFunc} can implement it, and what that implies about adapting plain functions to richer interfaces.
\end{exercise}

\begin{exercise}
Implement a stage between admission and projection that releases gated records after a named reviewer approves them. Decide which of the existing types should carry the approval, and whether the receipt of a record should change when it is approved.
\end{exercise}


\part{Concurrent Execution}
\chapter{Goroutines and Channels}
\label{ch:goroutines}

\section{What a Goroutine Is}

A goroutine is a function running concurrently with the rest of the program. It is started by placing the keyword \gocode{go} before a function call, and the call then proceeds on its own while the caller continues with the next statement. The Go runtime schedules goroutines onto a small number of operating-system threads, so they are cheap enough that a program can hold thousands of them, and a goroutine begins with a stack of only a few kilobytes that grows as needed. The cheapness changes how programs are designed, because the cost that makes a thread-per-task design impractical in many languages is largely absent, and a task that is naturally independent can simply be given a goroutine.

Three properties of goroutines are used throughout Part~III and should be stated at the outset. First, all the goroutines of a process share one address space, which means they can read and write the same variables, and the difficulty of concurrent programming is the difficulty of doing so safely. Second, when the \gocode{main} function returns, the process exits and every other goroutine is abandoned wherever it happened to be, without any chance to finish; a program that wants a goroutine's work must wait for it. Third, a goroutine has no identity visible to the program and no handle by which it can be stopped from outside, so a goroutine ends only when its function returns, and a program that wants to stop one must arrange for the function to notice and return. The second and third properties explain why the tools of this chapter, the wait group and the channel, are about coordination and not about starting work, which is the easy part.

\section{Why \golab{} Needs Them}

Verifying a file means reading its bytes and hashing them. For a handful of small files the work is trivial, and for a research directory holding thousands of large measurements it is the dominant cost of the \gocode{verify} command. The work is also independent from file to file: the digest of one file does not depend on any other, so nothing in the problem requires the files to be handled in sequence. The sequential version leaves most of the machine idle. While one file is being read the processor waits for the disk, and while its bytes are being hashed the disk waits for the processor, and several files in flight at once keep both busy.

The number of files in flight needs a bound, and the bound is a design decision and not an optimization detail. Every open file uses a descriptor, which the operating system limits, a disk serves many simultaneous readers worse than a few, and an unbounded program that started a goroutine per file would try to open every file in a large directory at the same moment. The standard remedy is a pool, a fixed number of worker goroutines that take jobs from a shared source until the source is exhausted. The package \gocode{pool} implements one for \golab{}.

\golisting{internal/pool/pool.go}{A bounded worker pool that verifies files concurrently.}{lst:pool}

\section{Reading the Pool}

The function \gocode{VerifyAll} makes three decisions about structure, each of which is a decision about who may touch what. It allocates the slice of results before any goroutine starts, with one element per input file. It creates an unbuffered channel \gocode{jobs} of integers, which will carry indices into the input. And it starts the workers, each a goroutine running the same closure, which receives indices from the channel until it is closed, processes the file at that index, and writes its outcome into the slot of the results slice with the same index.

A channel is a typed conduit through which goroutines pass values. A send on an unbuffered channel blocks until some other goroutine is ready to receive, and a receive blocks until some other goroutine sends, so the pair of operations is both a transfer of a value and a rendezvous between the two goroutines. The loop \gocode{for i := range jobs} receives repeatedly and ends when the channel has been closed and drained. Closing is the sender's business, and it is a signal to the receivers that no more values will arrive, so \gocode{VerifyAll} closes \gocode{jobs} after sending the last index. The rule that only the sender closes a channel is more than a convention, because sending on a closed channel causes a panic, and a receiver that closed a channel would be exposing every sender to it.

A \gocode{sync.WaitGroup} counts goroutines still running. The function calls \gocode{wg.Add(1)} before starting each worker and each worker calls \gocode{wg.Done()} in a deferred call when it returns, and \gocode{wg.Wait()} blocks until the count returns to zero. Placing the call to \gocode{Add} in the loop that starts the goroutines, and not inside the goroutines themselves, matters: if a goroutine called \gocode{Add} on its own behalf, the count could be zero when \gocode{Wait} was reached simply because the goroutine had not yet run.

The pool avoids two classical hazards by construction. It does not use a lock to protect the results, because no two goroutines ever touch the same element: each index is sent on the channel once and received by exactly one worker, which is therefore the sole writer of that slot. And it does not capture a loop variable in a goroutine closure. Before Go 1.22 a variable declared in a \gocode{for} clause was shared by all iterations, so a goroutine started in the loop that referred to it would see whatever value it held when the goroutine eventually ran. The module declares Go 1.21 and so has the older rule, and the code is written to give the same result under either rule, by sending the index through a channel and not by capturing it.

\section{Why the Caller Can Read the Results}

Writing to a slice from several goroutines and reading it from another raises the question of what the reader is guaranteed to see. The Go memory model answers it with a small set of ordering rules. A send on a channel happens before the corresponding receive completes, the closing of a channel happens before a receive that returns because the channel is closed, and a call to \gocode{Done} on a wait group happens before a \gocode{Wait} that it releases returns. An operation that happens before another is guaranteed to be visible to it, and two accesses to the same memory with no such ordering between them, at least one of which is a write, constitute a data race, whose effect the language does not define.

\begin{proposition}\label{prop:pool}
When \gocode{VerifyAll} returns, the result slice has one element per input file, element $i$ is the outcome for file $i$, and the caller's reads of the slice are ordered after every worker's writes to it.
\end{proposition}

\begin{proof}
The feeder sends each index $i$ exactly once, so exactly one worker receives it and writes element $i$, and no other goroutine accesses that element while workers are running. Every worker's last action is its deferred \gocode{Done}, and \gocode{VerifyAll} returns only after \gocode{Wait} has returned, so each write to an element happens before \gocode{Done}, which happens before \gocode{Wait} returns, which happens before the caller's reads. Because the feeder sends one index per file and receives block until sent to, every element is written before the function returns.
\end{proof}

The proposition explains what would go wrong if the call to \gocode{Wait} were removed. The function would return while workers were still writing, and the caller's reads of the slice would have no ordering with those writes, which is a data race even if the program happened to produce the right output on a given run.

\section{A Race, Observed}

The danger is easy to demonstrate and worth seeing once. The program in \texttt{examples/\allowbreak ch10\_\allowbreak race} starts a thousand goroutines, each of which increments a shared counter, and prints the total, which ought to be one thousand.

\golisting{examples/ch10_race/main.go}{A deliberately racy program.}{lst:race}

\lstinputlisting[style=shellstyle,caption={The racy program, run five times and then under the race detector.},label={lst:session3}]{session3.txt}

The five runs in Listing~\ref{lst:session3} give five different answers, and only the first is correct. The statement \gocode{counter++} is three steps, a read, an addition, and a write, and when two goroutines interleave their steps, one of the increments is lost. The result is not wrong in a consistent way that a test could pin down, and a program with this defect can pass a test suite for months. The race detector, enabled with the \gocode{-race} flag, instruments every memory access and reports a pair of accesses that are not ordered by the memory model, naming the line, the goroutines involved, and where they were created. A program in which it reports a race exits with status 66.

The detector is a tool and not a proof, and the book relies on it only in that sense. It reports races that occur in the executions it observes, so a race on a path the tests never exercise is invisible to it, and a program that passes under \gocode{-race} has been shown free of races only on the runs that were made. The tests of Part~III are therefore designed to drive concurrent paths hard, with many goroutines and many iterations, and are run with the detector enabled, and the arguments in the text, like the one above, are what establish that the code is correct for all runs.

\section{Goroutines Are Not Services}

It is tempting to treat a goroutine as a small independent service, a worker that is handed a job and returns a result, and the vocabulary of the pool encourages it. The comparison conceals the differences that matter most. Goroutines share memory and share fate: a panic in one goroutine that is not recovered terminates the entire process, a runaway goroutine slows every other, and a bug that corrupts a shared structure corrupts it for all of them. A service in the sense used in Part~V is a separately deployed program that communicates over a network, and its failure modes are of a different kind. A message can be lost or delivered twice, the remote program can be slow or absent or running an older version, and a request can be processed even though the reply never arrives, so that the sender cannot tell whether the work was done. None of these exists between goroutines, and designs that move from the first to the second without confronting them inherit failures that local concurrency never exposed. The pool of this chapter is a technique for using the processors and disks of one machine, and it makes no claim about distribution.

\golisting{internal/pool/pool_test.go}{Tests for the worker pool.}{lst:pool-test}

\section*{Exercises}

\begin{exercise}
Remove the call to \gocode{wg.Wait()} from \gocode{VerifyAll} and run the tests under \gocode{-race}. Report what the detector says and explain it in terms of Proposition~\ref{prop:pool}.
\end{exercise}

\begin{exercise}
The test \gocode{TestWorkersReallyRunInParallel} does not measure time. Explain how it proves that the workers run simultaneously, and why a test that compared elapsed times would be less reliable.
\end{exercise}

\begin{exercise}
Change the pool so that the workers pull file names from a shared slice protected by a mutex instead of receiving indices from a channel. Compare the two designs for clarity and for the guarantees available to the caller.
\end{exercise}

\begin{exercise}
Write a program that starts a goroutine which panics and show that the process terminates even though \gocode{main} is waiting normally. Then add a \gocode{recover} in the goroutine and discuss whether recovering is appropriate for a verification worker.
\end{exercise}

\chapter{Contexts, Select, and Cancellation}
\label{ch:contexts}

\section{The Problem of Stopping}

Starting concurrent work is easy and stopping it is hard. A goroutine cannot be killed from outside, as Chapter~\ref{ch:goroutines} noted, so any work that may need to be abandoned must be written to notice a request to stop and to return of its own accord. The request has to travel from whoever wants the work stopped to whoever is doing it, and the path may pass through several layers of functions and several goroutines, none of which knows about the others. The standard library solves the problem with a single type, \gocode{context.Context}, which carries a cancellation signal and an optional deadline down through a program, and which has become the first parameter of nearly every function in Go that can block.

A context is created from another. The function \gocode{context.Background} returns an empty context that is never cancelled and serves as the root. The function \gocode{context.WithCancel} derives from a parent a child context together with a function that cancels it, and \gocode{context.WithTimeout} derives a child that cancels itself when a duration has elapsed. Cancellation flows downward: cancelling a context cancels every context derived from it, and cancelling a child leaves the parent untouched. The tree of contexts therefore mirrors the tree of work, and cancelling a branch stops the work beneath it and nothing else.

Two methods are used to observe a context. The method \gocode{Done} returns a channel that is closed when the context is cancelled, which makes cancellation something a goroutine can wait for with the \gocode{select} statement described below. The method \gocode{Err} returns \gocode{nil} while the context is live and afterwards returns \gocode{context.Canceled} or \gocode{context.DeadlineExceeded}, saying why it ended, and it is the cheaper form when the code only needs to ask whether to continue. The runner of Chapter~\ref{ch:processes} already used a context with a timeout to bound a child process, and the pool applies the same idea to a whole batch of files.

\section{The Select Statement}

A goroutine often needs to wait for more than one thing, such as a result, a cancellation, or a timer, and to act on whichever happens first. The \gocode{select} statement does this. It lists several channel operations, blocks until one of them can proceed, and runs the corresponding branch. If several can proceed at the same moment, the runtime picks one at random, so a program must not depend on an order among simultaneously ready cases, and a \gocode{default} branch, if present, runs immediately when none is ready, which turns the statement into a non-blocking poll.

The tests of the pool use \gocode{select} where a bounded wait is needed. In \texttt{Test\-Workers\-Really\-Run\-In\-Parallel}, the test waits for each of four workers to announce that it has arrived at the verification function, and each wait is a \gocode{select} between the arrival channel and a timer created by \gocode{time.After}. If the workers were serialized, the second arrival would never come, the timer would fire after five seconds, and the test would fail with a message that says how many workers were running. A test that waited on the channel alone would hang forever in that case, and a hung test is a poor way to report a defect.

\section{The Pool's Cancellation Contract}

The pool takes a context as its first argument, in the usual position, and its documentation states what cancellation does. Files that no worker has started when the context ends are not verified, and each of them receives the context's error as its outcome. A verification already under way is not interrupted. Two decisions are contained in that sentence.

The first is that cancellation does not shorten the result. The feeder loop continues to send every index even after the context ends, and the workers, seeing the cancelled context, record the cancellation as the outcome of each remaining file instead of verifying it. The alternative design would stop feeding at the first sign of cancellation and return a short slice, or a slice with empty elements for the files never reached. Either would leave the caller to work out which files had been skipped, and for a tool whose purpose is an honest record, a silent omission is the worst kind of result. The choice is the one made in Chapter~\ref{ch:errors} for partial failure: every file ends up with an outcome, and the outcome says what happened, which in this case is that it was not attempted because time ran out.

The second decision is that a running verification is not interrupted, because \gocode{artifact.Verify} takes no context and reads the file to the end. For ordinary files the delay is small, and for a very large file on a slow disk it is not. The honest statement is that the timeout bounds when new work is started and does not bound when the call returns, and the documentation says so. Making verification itself cancellable is straightforward in principle, since the copy loop can check the context between blocks, and it is the first exercise of this chapter.

\golisting{internal/pool/pool_test.go}{Tests for the worker pool, repeated for reference.}{lst:pool-test-again}

Two tests pin down the contract. \texttt{Test\-Cancelled\-Context\-Still\-Reports\-Every\-File} cancels the context before the call and checks that all twenty files come back with \gocode{context.Canceled} and that the verification function was never invoked. \gocode{TestCancellationMidRun} is more delicate. It uses one worker and a verification function whose first call cancels the context through a \gocode{sync.Once}, and it checks that exactly one file was verified and the other forty-nine were cancelled. The test is deterministic only because it uses one worker. With several workers, the number of files already in progress when the cancellation lands varies from run to run, and a test that asserted an exact count would be flaky. The contract promises only that every file has an outcome and that files not started are cancelled, and the tests assert precisely that and no more.

\section{Cancellation in the Command}

The \gocode{verify} command exposes the mechanism through the flag \gocode{-timeout}. When it is given, the command derives a context with that timeout and passes it to the pool. The flag \gocode{-workers} sets the size of the pool, defaulting to the number of processors, and the flag \gocode{-each} asks for one ledger record per file, which is the subject of the next chapter.

\lstinputlisting[style=shellstyle,caption={Concurrent verification, per-file records, and a timeout that expires at once.},label={lst:session5}]{session5.txt}

The transcript in Listing~\ref{lst:session5} shows all three effects. The first command verifies six files with four workers and writes a record for each file followed by the aggregate record, seven in all, and the audit confirms that the chain holds. The order in which the per-file records appear is not the order of the files: \texttt{d04} precedes \texttt{d03} and \texttt{d02}, because the ledger records files in the order the workers finished them, and the order of completion is a property of the scheduler that the program does not control. The printed digests, by contrast, are in input order, because the command prints from the result slice after the pool has returned. The two orders serve different readers, and the difference is intended. The ledger is a history of what happened, and the terminal output is a report in a stable order.

The second command sets a timeout of one nanosecond, which expires before any work can start. All six files come back with the deadline error, the command exits with status 1, and the aggregate record in the ledger states that six files failed. The attempt is on the record. An operator who looks at the ledger a week later sees that a verification was requested, that nothing was verified, and why, which is the same principle that governed the ordering of the \gocode{verify} command in Chapter~\ref{ch:errors}.

\section*{Exercises}

\begin{exercise}
Change \gocode{artifact.Verify} to accept a context and to check it between blocks of a fixed size while hashing. Choose the block size and justify it, and write a test that cancels a verification of a large file part-way through. Decide what the returned error should say about the bytes already read.
\end{exercise}

\begin{exercise}
Write a version of the feeder that stops sending as soon as the context is cancelled and fills the unsent slots with the context's error afterwards. Show that it satisfies the same contract, and identify which design is easier to reason about.
\end{exercise}

\begin{exercise}
Replace the use of \gocode{ctx.Err()} in the worker with a \gocode{select} on \gocode{ctx.Done()} and a \gocode{default} case. Explain whether the two behave differently, and why one might prefer either.
\end{exercise}

\begin{exercise}
Make the workers' error for a cancelled file wrap \gocode{context.Canceled} or \gocode{context.DeadlineExceeded} with the path of the file, and show that \gocode{errors.Is} still recognizes the cause.
\end{exercise}

\chapter{Shared State and the Single Writer}
\label{ch:writer}

\section{The Fork That Part I Predicted}

Chapter~\ref{ch:ledger} ended with a warning. The ledger's \gocode{Append} reads the whole file, checks the chain, computes the next record, and writes it, and these steps are not one atomic action. Two processes that run them at about the same time can each read the same last record and each write a successor, so that the chain acquires two records with the same sequence number. The warning can be tested by launching twenty copies of the command at once, each of which runs a trivial program and records the outcome.

\lstinputlisting[style=shellstyle,caption={Twenty processes appending to one ledger.},label={lst:session4}]{session4.txt}

The result in Listing~\ref{lst:session4} is the failure predicted. Of twenty commands, five managed to append, and two of them wrote sequence number four. The audit reports the damage at the fifth record. The remaining fifteen commands exited with an error, because by the time they ran, the ledger already failed its check and \gocode{Append} refuses to extend a broken chain. The behaviour is not an accident of the transcript. The exact numbers depend on timing and would differ on another run, but the shape is general: once two writers collide, the ledger is damaged, the damage is detected, and every later writer is turned away. The design failed closed, and the cost was an unusable ledger.

The defect is a race of the kind the previous chapters studied, with the difference that the shared state is a file and the competing parties are processes and not goroutines. The pool of Chapter~\ref{ch:goroutines} avoided shared writes by giving each slot of the result slice to one goroutine. The ledger cannot be partitioned that way, because every record depends on the one before it, and the whole point of the structure is a single total order. A total order needs exactly one place at which the order is decided.

\section{Two Ways to Serialize}

There are two standard ways to make a shared resource safe for many goroutines. The first is a lock: a \gocode{sync.Mutex} guards the resource, every user acquires it before touching the resource and releases it afterward, and the lock ensures that only one goroutine is inside the guarded section at a time. The second is ownership: a single goroutine is given exclusive responsibility for the resource, and every other goroutine asks it to act by sending a message over a channel. Go programmers summarize the second style with a proverb, that one should share memory by communicating and not communicate by sharing memory. The proverb is a preference and not a rule, and either style can be right. A lock is the simpler tool for protecting a small piece of state with a few operations. Ownership suits a resource that has a lifecycle, an order of operations, and a single point where work should be sequenced, and the ledger has all three.

\section{The Writer}

The type \gocode{ledger.Writer} implements ownership. Its constructor starts one goroutine, the loop, which receives requests from a channel and calls the existing \gocode{Append} function for each, one at a time, so that within the process no two appends ever overlap. Every other goroutine calls the method \gocode{Append} on the writer, which packages the arguments into a request, sends it, and waits for a reply. The listing also contains a method \gocode{Snapshot}, which sends a request of a second kind that asks the loop to read the whole ledger instead of appending to it. It is not needed in this chapter and is explained in Chapter~\ref{ch:http}, where a service has readers that run while appends are in progress.

\golisting{internal/ledger/writer.go}{The single-writer ledger.}{lst:writer}

Several details of Listing~\ref{lst:writer} carry weight. Each request contains its own reply channel with a buffer of one, so the loop can deliver the result without waiting for the caller to be ready, and a caller that has gone away cannot block the loop or leave a goroutine stranded. The detail map is copied before it is queued, for the reason given in Chapter~\ref{ch:interfaces}: a map is a reference, and a caller that reused its map for the next call would otherwise change a request the writer had not yet processed. The timestamp is taken inside the loop, at the moment of writing, and not by the caller, so the times in the ledger never run backward relative to the sequence numbers.

The subtle part is shutdown. A channel can be closed to tell its receivers that no more requests will come, and closing it ends the loop, but sending on a closed channel panics, so closing the channel while another goroutine is about to send would crash the program. The writer prevents this with a read-write lock. Each call to \gocode{Append} holds the read lock while it checks the closed flag and sends its request, and \gocode{Close} takes the write lock to set the flag and close the channel. The write lock waits until all current holders of the read lock have released it, which means all in-flight sends have completed, and the loop is still receiving during that time, so they do complete. Once \gocode{Close} holds the lock and has closed the channel, any later \gocode{Append} sees the flag and returns \gocode{ErrClosed} without sending. The loop drains whatever is already in the channel before it ends, so every request accepted before the close is written.

\begin{proposition}\label{prop:writer}
Suppose a single \gocode{Writer} is the only code appending to its file. Then every call to \gocode{Writer.Append} that returns a nil error corresponds to exactly one record on disk, the records appear in the order in which the loop received the requests, and the chain of digests remains valid. After \gocode{Close} returns, no further record is written.
\end{proposition}

\begin{proof}
The loop is the only goroutine that calls \gocode{ledger.Append} on the file, and it handles one request completely before receiving the next, so appends are sequential and each reads the file as left by the one before. Each is the same function whose correctness for sequential use was established by the tests of Chapter~\ref{ch:ledger}. A request is replied to after its append returns, so success reported to a caller means the record was written. The channel is closed only under the write lock, after which the closed flag prevents any send, and the loop ends only after the channel is drained, so requests accepted before the close are processed and requests after it are refused.
\end{proof}

The proposition has a hypothesis, and the hypothesis is the limit of the design. The writer protects the file from goroutines in one process. It does nothing about a second process, which has no knowledge of the first writer's goroutine and will read and write the file on its own account. Twenty copies of the command, each with its own writer, would reproduce the failure of Listing~\ref{lst:session4} exactly. Protecting the file against other processes needs a mechanism from the operating system, such as an advisory lock on the file, or an arrangement in which only one process ever writes, and the second of those is the shape that Part~V gives the tool by making a service the sole owner of the ledger and having clients send it requests over a network. The first exercise explores the lock.

\golisting{internal/ledger/writer_test.go}{Tests for the single-writer ledger.}{lst:writer-test}

\section{What the Tests Establish}

The first test starts twenty goroutines that each append ten records through one writer, then checks that all two hundred records are present, that the chain holds, and that every subject appears exactly once, so that no record was lost and none was written twice. The test was run with the race detector enabled. The third test is the one that exercises the shutdown protocol. It starts sixteen goroutines appending twenty records each, calls \gocode{Close} from another goroutine at an arbitrary moment, and checks that the number of appends that reported success equals the number of records on disk and that the chain is valid. Every append is either accepted and written or refused with \gocode{ErrClosed}, and none is lost in between. As with every concurrency test in the book, the test samples some interleavings and does not enumerate them, and the argument that the protocol is correct is the one in Proposition~\ref{prop:writer}.

With a writer in place the command can offer what Chapter~\ref{ch:errors} said it lacked. The flag \gocode{-each} gives the \gocode{OnResult} callback of the pool a function that appends a record for each file as soon as the file's outcome is known. The callbacks run in the worker goroutines, so up to as many appends may be requested at once as there are workers, and the writer serializes them. The transcript in Chapter~\ref{ch:contexts} shows the resulting records, in order of completion and with a valid chain, followed by the aggregate record.

\section{Costs and Where Part III Ends}

The writer is correct and not fast. Each call to the underlying \gocode{Append} reads and checks the entire ledger, so the cost of the $n$th append is proportional to $n$ and the cost of $n$ appends is proportional to $n^2$. For a ledger of a few thousand records the delay is imperceptible, and for a ledger of a million it is prohibitive. A writer that holds the head digest and the next sequence number in memory, and checks the whole chain only when it opens the file, would make each append constant in cost, at the price of trusting its own memory between appends. The trade is a common one in systems that keep an authoritative log, and the third exercise asks for it.

Part~III has given the tool concurrency in three forms, a bounded pool for independent work, a context for cancellation, and a single writer for the one resource that must be sequenced. Each is accompanied by a statement of what it does not do. The pool does not interrupt a running verification, the context bounds the start of work and not its end, and the writer does not protect against other processes. Part~IV turns to data that must outlive a process, and the question there will be which of what \golab{} knows should be stored, and which can always be rebuilt.

\section*{Exercises}

\begin{exercise}
Add an advisory file lock around the body of \gocode{ledger.Append}, using the facility your operating system provides, and rerun the twenty-process experiment. Report what changes, and identify one situation in which the lock fails to protect the chain, such as a ledger on a network filesystem.
\end{exercise}

\begin{exercise}
Replace the single-writer goroutine with a \gocode{sync.Mutex} around \gocode{ledger.Append}, and compare the two designs with respect to shutdown, to the order of records, and to the amount of code. State which you would choose and why.
\end{exercise}

\begin{exercise}
Make the writer keep the head digest and the next sequence number in memory, checking the chain on disk only when the writer is created. Show that appending a thousand records is faster, and describe a scenario in which the in-memory state and the file disagree.
\end{exercise}

\begin{exercise}
Remove the read lock from \gocode{Writer.Append} and write a test that calls \gocode{Close} while appenders are active, run under the race detector and repeated many times. Describe the failure that appears and relate it to the rule about closing channels.
\end{exercise}


\part{Memory and Reconstruction}
\chapter{Encoding, Buffering, and Atomic Replacement}
\label{ch:encoding}

\section{Data That Outlives a Process}

Everything built so far either finishes within one run of the program or goes into the ledger, which is written once and never rewritten. Part~IV is concerned with a third kind of data, which is written repeatedly, replaced as a whole, and read by later runs of the program or by other programs. A file of this kind raises a set of questions that the ledger avoided by being append-only. How is a structure turned into bytes and back without losing or inventing anything? What happens when the bytes are large? What does a reader see if it opens the file while a writer is part-way through replacing it, and what is left behind if the machine stops in the middle? Go answers the first question with the \gocode{encoding} packages, the second with \gocode{bufio}, and the third with no package at all, only a technique that depends on a property of the operating system, which this chapter states exactly and then tests as far as a test can.

\section{Encoding as a Round Trip}

The ledger already uses \gocode{encoding/json}, and it is worth saying what the package does with a struct, because the later chapters depend on it. The encoder writes the exported fields of a struct in the order of their declaration, uses the name in the field tag when there is one, and omits a field marked \gocode{omitempty} when its value is the zero value. It writes the keys of a map in sorted order. These rules make the encoding of a value a function of the value alone, which is why the digests of Chapter~\ref{ch:ledger} are reproducible. The decoder is more permissive than the encoder. It ignores fields it does not recognize, leaves absent fields at their zero values, and accepts the fields in any order, so a file written by a newer version of a program is read without complaint by an older one, and the fields the older version does not know about are silently dropped when it writes the file again.

That permissiveness is convenient and also a hazard, and the index of the next chapter guards against it with an explicit version number and with checks on the decoded value, for instance that no path appears twice. The general point is that a round trip through an encoding preserves exactly what the struct can express. Information that the struct has no field for does not survive, and the compiler has nothing to say about it. A reader who wants a round trip to be faithful has two honest options, which are to decode into a type that rejects unknown fields by calling \gocode{DisallowUnknownFields} on the decoder, or to state in the documentation that unknown fields are discarded. The index takes the second course, because an index is derived data and a field it fails to preserve can be recomputed.

\section{Buffering and a Limit That Was There All Along}

Reading and writing a file one small piece at a time is slow, because each call to the operating system has a fixed cost. The package \gocode{bufio} wraps a reader or writer in a buffer so that many small operations become few large ones. A \gocode{bufio.Writer} holds its data until the buffer fills or until \gocode{Flush} is called, and the last of these is the one that is easy to forget: bytes written to a buffered writer and not flushed are lost without any error. A \gocode{bufio.Scanner} reads a file in lines, tokens, or words, and the ledger uses it to read one JSON record per line.

The scanner has a limit that the ledger of Part~I did not mention. A scanner holds a line in a buffer of bounded size, and when a line is longer than the limit it stops with the error \gocode{bufio.ErrTooLong}. The ledger sets the limit at one mebibyte, which is far more than any record the program produced in the earlier parts. The limit is nonetheless a property of the reader that the writer did not respect. \gocode{Append} would write a record of any length, and a record longer than the limit could therefore be written successfully and then never read back, and since \gocode{Append} begins by loading the whole ledger, every later append would fail as well. A single oversized detail value, such as a long error message, would disable the ledger permanently.

The repair is small and the lesson is general. The ledger package now declares the limit as an exported constant, \gocode{MaxLine}, sets the scanner's buffer from it, and refuses in \gocode{Append} to write a record whose encoding exceeds it, returning the sentinel \gocode{ErrTooLarge}. The new test appends a record that is too large, checks that the error is the sentinel, and then checks that a normal append still succeeds and that the ledger loads. The principle is that a writer must not produce what its own reader cannot consume. Where the format has a limit, the limit belongs in one place, and the code that writes must check it as well as the code that reads.

\section{Replacing a File Whole}

The ledger can be appended to because each record is independent of the file's other contents. A file that is rewritten as a whole, such as the index, is in a different position. The obvious way to rewrite it is to open the file for writing, which truncates it, and write the new contents. That sequence has a window in which the file is empty or partly written. A reader that opens the file in that window sees a fragment, and a crash in that window leaves a fragment behind permanently, with the old contents already destroyed.

The standard remedy is to write the new contents to a different file and then rename it over the old one. On POSIX filesystems, renaming a file to a name that already exists, within one filesystem, replaces the destination in a single step, so that any process that opens the name sees either the entire old file or the entire new one. The package \gocode{atomicfile} packages the technique.

\golisting{internal/atomicfile/atomicfile.go}{Replacing a file by writing a temporary and renaming it.}{lst:atomicfile}

The function creates the temporary file in the same directory as the destination, because a rename between directories on different filesystems is a copy and has no such guarantee. It hands the caller a buffered writer through a callback, so that the caller can stream a large structure without first building all its bytes in memory, and it flushes the buffer itself, so that the caller cannot forget. It sets the permissions, calls \gocode{Sync} on the file to ask the operating system to write the data to stable storage before the rename, closes the file and checks the error, renames, and finally syncs the directory, because the rename is itself a change to the directory and is durable only once the directory has been flushed.

The cleanup deserves a comment because it uses a feature of the language that is easy to misread. The function has a named result, \gocode{err}, and a deferred function that inspects it. Every failing step assigns to that variable before returning, so the deferred function can tell whether the function is failing and, if so, close the temporary file when it has not been closed and remove it. A deferred function runs after the return value has been set, and with a named result it can read what was set, which is how a single piece of cleanup code covers every early return. The variable \gocode{closed} records that the file was already closed, since closing a file twice reports an error that would mask the real one.

\begin{proposition}\label{prop:atomic}
Suppose \gocode{Write} returns a non-nil error. Then the file at the destination path has the same contents it had before the call, and no temporary file created by the call remains, except possibly when the error comes from the final synchronization of the directory, in which case the new contents are in place.
\end{proposition}

\begin{proof}
The destination is modified only by the call to \gocode{os.Rename}, which comes after every step that can fail except the directory synchronization. If any earlier step fails, the destination has not been touched, and the deferred function removes the temporary file. If the rename itself fails, the destination is unchanged because a rename either takes effect or does not, and the deferred function removes the temporary file. If only the directory synchronization fails, the rename has taken effect, the temporary name no longer exists, and the error reports that the replacement may not survive a crash.
\end{proof}

The proposition is about the function's behaviour on errors it can observe. It does not cover a crash, since a program that has stopped cannot clean up. A crash before the rename leaves a stray temporary file with a distinctive name and the destination untouched. A crash after the rename but before the directory flush may, on some filesystems, leave the old file in place after the machine restarts. Neither case produces a partial destination file, which is the property the technique exists to provide, and neither can be demonstrated by an ordinary test, because a test that is running has not crashed.

\section{What the Tests Establish}

\golisting{internal/atomicfile/atomicfile_test.go}{Tests for atomic replacement.}{lst:atomicfile-test}

The tests cover what an ordinary test can reach. The first checks that a file is created and then replaced, that the permissions are applied, and that no extra files remain in the directory. The second makes the callback write more bytes than the buffer holds, so that some of them reach the temporary file, and then return an error. It checks that the old contents are intact and that nothing is left over, which tests the first sentence of the proposition at the point where it is most at risk. The third checks that a missing directory is reported as an error. The fourth runs a reader in a loop while the main goroutine replaces a one-mebibyte file forty times, alternating between two contents that consist entirely of one repeated letter, and checks that every read returns one of the two complete contents. The test samples the interleavings of one run on one filesystem and does not prove the property for every filesystem, and the reader of the chapter should treat the rename guarantee as one that the operating system makes and that the test merely fails to contradict.

\section*{Exercises}

\begin{exercise}
Change the ledger's \gocode{Load} to use a \gocode{json.Decoder} reading directly from the file, and not a scanner, so that no line limit exists. Decide whether the constant \gocode{MaxLine} and the error \gocode{ErrTooLarge} should then be removed, and justify the decision in terms of what a very large record would do to the memory use of \gocode{Append}.
\end{exercise}

\begin{exercise}
Remove the call to \gocode{Sync} on the temporary file and describe a failure that the tests of this chapter cannot detect. Explain why the order of the synchronization and the rename matters.
\end{exercise}

\begin{exercise}
Make the temporary file in a different directory from the destination, on a different filesystem if one is available, and report what the rename does. State what this implies about the correct location of the temporary file.
\end{exercise}

\begin{exercise}
Modify the decoding of the index so that unknown fields are an error, using \gocode{DisallowUnknownFields}. Describe a situation in which this makes an upgrade of the program harder, and a situation in which it catches a real mistake.
\end{exercise}

\chapter{Derived Data and the Index}
\label{ch:index}

\section{Authority and Convenience}

The ledger answers questions slowly. To learn the last digest recorded for a given file, a program must read every record from the start and keep the last one that mentions the file. That is acceptable for a ledger of a few hundred records and tiresome for a ledger of a million. The usual response is to keep a second structure, an index, that holds the answer to the common questions directly. The response is safe or unsafe according to a single decision, which is whether the program treats the index as a source of truth or as a summary of one. In the first case the index must be protected, backed up, and kept consistent with everything else, and a loss of the index is a loss of information. In the second case the index is derived data. It may be deleted at any time, and the only cost is the time to compute it again.

The index of \golab{} is of the second kind, and the design is arranged to make that status checkable and not merely declared. Three rules are enforced in the code. The index is computed from the ledger records alone and from nothing else. It records which prefix of the ledger it was computed from, by the number of records and the digest of the last one. And the program never makes a decision from the index that it could not make by recomputing the index from the ledger, a rule that holds because every path that uses the index first obtains it from a function that falls back on rebuilding.

\section{The Package}

\golisting{internal/index/index.go}{The index, derived from the ledger.}{lst:index}

An \gocode{Entry} summarizes what the ledger says about one path: the digest from its most recent successful verification, whether the latest attempt succeeded, the number of attempts and failures, and the sequence number of the last record that mentioned it. The fields are chosen for the questions the command will ask, and the entry deliberately keeps the last known digest even when a later attempt failed, so that a path whose verification failed once is not mistaken for a path that was never verified. The \gocode{Index} itself holds the entries in a map from path to entry, together with the number of records consumed and the digest of the last. The map is unexported, so that the only ways to change an index are the methods of the package, and every change goes through \gocode{Apply}.

The method \gocode{Apply} consumes one ledger record and begins by checking that the record is the immediate successor of the last one consumed: its sequence number must be one more than the count, and its predecessor field must equal the remembered head. A record that skips ahead, repeats, or belongs to a different chain is refused with \gocode{ErrOutOfSequence}. This check is what makes the pair of numbers a faithful description of the index. The count and the head are not labels attached from outside; they are advanced by the same method that updates the entries, and the method will not advance them over a gap. Only the two verification actions affect entries. Every other record, such as a run or an admission, advances the position and leaves the entries alone, so the index covers the whole ledger even though it summarizes only part of what is in it.

The function \gocode{Build} computes the index of a complete ledger by applying every record to an empty index, and it begins by calling the ledger's \gocode{Check}. A table derived from a broken chain would look authoritative while resting on nothing, and the function therefore refuses to produce one. The method \gocode{CatchUp} brings an existing index forward. It verifies that the index belongs to the ledger it is given, by checking that the record at the index's position has the digest the index remembers, then checks the chain and applies the records that follow. A shorter ledger, or a ledger whose record at that position has a different digest, produces \gocode{ErrStale}.

The function \gocode{Open} ties the pieces together. It reads the index file if there is one, decodes it, and tries to catch it up. If the file is missing, unreadable, of another version, inconsistent, or not a prefix of this ledger, it builds the index from the ledger instead, and it returns the reason. The result says which of the three things happened, so that a command can tell the operator whether the index was loaded as it stood, caught up with new records, or rebuilt. The one case in which \gocode{Open} does not fall back is a ledger that fails its own check, because there is nothing authoritative left to rebuild from, and the error is returned as it is.

\section{Reconstruction Equals Incremental Update}

An index that can be updated by appending records and also rebuilt from scratch raises an obvious question, which is whether the two routes give the same result. If they could differ, the index would depend on its history, and the claim that it is derived from the ledger alone would be false.

\begin{proposition}\label{prop:rebuild}
Let $L$ be a ledger of $n$ records that satisfies \gocode{Check}, and let $0 \le k \le n$. The encoding of the index obtained by building from the first $k$ records and then calling \gocode{CatchUp} with $L$ is byte for byte the encoding of the index built from all of $L$.
\end{proposition}

\begin{proof}
\gocode{Build} is the fold of \gocode{Apply} over the records from the empty index, and \gocode{CatchUp} continues the same fold from the position $k$ over the remaining records. The state of an index is its count, its head, and its map of entries, and \gocode{Apply} computes the new state from the old state and the record by a deterministic rule that consults nothing else, no clock, no file, and no iteration over the map. After $n$ applications the two routes have applied the same rule to the same records in the same order starting from equal states, so they end in equal states. The encoding of an index depends only on its state, because the entries are written sorted by path and the count and head are written as they stand. Equal states therefore have equal encodings.
\end{proof}

The proof rests on facts that the code could lose without the compiler noticing. If \gocode{Apply} read the current time, or consulted an environment variable, or built part of its state by ranging over a map and so depended on the randomized iteration order, the proof would break and the proposition would be false. The sorted output of \gocode{Entries} exists for the same reason, and the test that checks the proposition does so for every cut point $k$ from zero to the length of the ledger. As in the earlier chapters, the test checks instances and the argument covers the cases, and the argument is only as good as the discipline that keeps \gocode{Apply} pure.

\golisting{internal/index/index_test.go}{Tests for the index.}{lst:index-test}

\section{What the Index Does Not Protect Against}

The chain of digests gives the index a way to notice that it belongs to a different ledger or has fallen behind, and it gives no way to notice that someone edited the index. The head and the count are checked against the ledger. The entries are not, because the check that would compare them is the check that would rebuild them. An index file whose head and count match the ledger but whose entries have been altered is accepted by \gocode{Open} and used. The fault is not in the code, which does what it says. It follows from the structure: a summary that is cheaper than the source cannot be verified against the source without paying the cost of the source.

Two responses are available, and the book takes the first. The index is treated as untrusted input whose effect is bounded. The command that uses it, which the next chapter presents, always hashes the files on disk itself and uses the index only to remember what the ledger last recorded, so a falsified entry can produce a wrong report that a file changed or did not change, and it cannot produce a record in the ledger or a decision to admit anything. The command that writes the index offers a flag that rebuilds it unconditionally, which is the remedy when an operator has reason to doubt the file. The second response is to store in the index a digest of its own entries, which detects accidental damage and not a deliberate edit by someone who can recompute it, and the third exercise asks for it.

\section*{Exercises}

\begin{exercise}
Extend \gocode{Entry} with the time of the first and last attempts, taken from the ledger records. Explain why the time must come from the record and not from the clock at the moment of indexing, and show how the proof of Proposition~\ref{prop:rebuild} would fail if it came from the clock.
\end{exercise}

\begin{exercise}
Add a second index, from action name to count, and show that the single pass in \gocode{Build} can fill both. Discuss whether the two should be one type or two, and how \gocode{Open} would decide that either is stale.
\end{exercise}

\begin{exercise}
Store in the index file a digest of its entries and check it on load. State exactly which edits the check detects and which it does not, and explain why it does not make the index authoritative.
\end{exercise}

\begin{exercise}
The method \gocode{CatchUp} calls \gocode{Check} on the whole ledger on every call, which makes catching up proportional to the size of the ledger and not to the number of new records. Propose a variant that checks only the new records and the link to the old head, state what it assumes about the old records, and say when that assumption is unsafe.
\end{exercise}

\chapter{Reconstruction and Its Limits}
\label{ch:reconstruction}

\section{Two Commands}

Two commands put the index to use. The command \gocode{index} obtains an index that agrees with the ledger, by loading, catching up, or rebuilding as \gocode{Open} decides, writes it to disk with the atomic replacement of Chapter~\ref{ch:encoding}, and reports which of the three happened and why. The command \gocode{status} answers the question that motivated the index, which is whether the files in a directory still match what the ledger last recorded about them. It hashes every file under the directory with the pool of Part~III, looks each path up in the index, and prints one line per path with a state. A path is \emph{unchanged} when its current digest equals the last recorded digest and the latest attempt succeeded. It is \emph{changed} when the digests differ, \emph{new} when the ledger has never mentioned it, \emph{missing} when the ledger mentions it under the directory and no file is there now, and \emph{unreadable} when it could not be hashed. Two further states, \emph{never-verified} and \emph{last-attempt-failed}, report paths whose history in the ledger contains only failures or ends in one. The command exits with a failure status when any path is not unchanged, so that a script can use it as a check, and it writes nothing at all: not to the ledger, not to the index.

\section{A Recorded Session}

The session below uses a directory of three small files. The first command verifies them with one worker and with the \gocode{-each} flag, so that the ledger receives a record for each file in a deterministic order. The rest of the session builds the index, destroys it, builds it again, and then disturbs the files and the ledger in turn.

\lstinputlisting[style=shellstyle,caption={Building, destroying, and rebuilding the index, and the states reported by \gocode{status}.},label={lst:session6}]{session6.txt}

The transcript carries the argument of Part~IV in a few observations. The first build says that it rebuilt the index because no index file existed, and that it covers three paths from four ledger records, the three file records and the aggregate. The listing of the start of the file shows the version, the count of records, the head digest, and the first entry, in an encoding whose layout is fixed by the declaration of the types. The checksum of the file is printed, the file is deleted, the command is run again, and the checksum is the same. That is Proposition~\ref{prop:rebuild} observed on one ledger: the file was a function of the ledger, and a function recomputes to the same value.

The session then changes the world. The file \texttt{b.txt} is overwritten, a new file \texttt{d.txt} is created, and \texttt{c.txt} is removed, and \gocode{status} reports a changed, a new, and a missing path, with the other still unchanged and an exit status of 1. The report says nothing about why the content of \texttt{b.txt} changed or whether the change was legitimate. It says that the bytes on disk no longer match the digest the ledger holds, which is a statement about a difference and not about its cause.

After the files are verified again, the ledger gains four more records, and the command \gocode{index} reports that it caught up, with four paths from eight records, so that the existing file was extended and not rebuilt. A new \gocode{status} shows \texttt{b.txt} and \texttt{d.txt} as unchanged and \texttt{c.txt} still as missing, and the last of these is correct and worth noticing. The ledger is a history, and it records that \texttt{c.txt} was once verified. Deleting the file does not delete the fact that it was verified, and the index, being derived from the ledger, keeps reporting the path until someone records otherwise. A tool that wanted removals to clear the entry would need a ledger action for removal, which is a design decision about the ledger and not a detail of the index.

The last two commands follow the alteration of one letter in the second line of the ledger. Both \gocode{status} and \gocode{index} stop with the same error, naming record 2 and the broken chain, and neither produces output from the damaged ledger. The index file that existed on disk was not consulted in a way that could have hidden the damage, because \gocode{Open} checks the chain before it trusts anything.

\section{What Was Not Shown}

Two properties of the design were checked separately and are reported here without a transcript. The output of \gocode{status} is the same whether or not the index file exists. In a directory prepared like the one above, running the command with the file present, deleting the file, and running it again produced identical output, which is what the rule that no decision depends on the index predicts. The second property is the safety of concurrent writers. Thirty simultaneous runs of \gocode{index} against one ledger left a single index file with the same checksum as a single run, and no temporary files behind. The result is the one the rename argument gives. Each process writes its own temporary file and renames it over the destination, every one of the files is a correct index of the same ledger, and so whichever rename lands last, the destination is correct. The contrast with the ledger of Chapter~\ref{ch:writer} is instructive. The ledger is a sequence, each record depends on the one before, and two writers disagree about what comes next. The index is a function of the ledger, any two writers compute the same value, and a collision between them costs nothing. Derived data is cheap to make safe for concurrency because it has no history of its own to protect.

\section{Limits of the Analogy}

The part is named for memory and reconstruction, and the name needs a statement of what it claims. The engineering discipline is plain: keep one authoritative record, derive everything else from it, and be able to delete the derived forms without loss. The word \emph{reconstruction} is used for that discipline and for nothing more. Nothing in this chapter suggests that an index resembles a memory in how it is formed, how it fails, or what it is for. A human memory is not a function of a log, and it changes when it is used. A cache is exactly a function of a log and does not change. The two share one property, that each is rebuilt from something other than itself, and the book makes no further use of the comparison.

The engineering limits are several. Rebuilding costs time proportional to the ledger, so an index is a saving only when it is loaded and extended more often than it is rebuilt. The derived data is only as good as the authoritative record, and the ledger records what the program was told and what it did, not what is true of the files in the world, as the persistent \emph{missing} entry shows. The index is trusted to the degree that its file has not been edited, which Chapter~\ref{ch:index} identified as an unchecked risk with a bounded effect. And the chain beneath everything is tamper-evident only relative to an external copy of its head, a limit that Chapter~\ref{ch:ledger} stated and that nothing in Part~IV removes.

Part~IV leaves the tool with a durable ledger, a derived index, and a safe way to replace a file. All of it runs inside one machine and is operated by one person at a command line. Part~V removes that restriction by making the tool a service that other programs can reach over a network, and the first thing the service will need is what the ledger could not have for itself in Part~III: a single owner for the one resource that must be sequenced.

\section*{Exercises}

\begin{exercise}
Add a ledger action that records the removal of a path, have \gocode{status} report a removal as a distinct state, and update \gocode{Apply} so that an entry is cleared by a removal record. Decide whether the index should drop the entry or keep it marked as removed, and justify the choice in terms of the proposition on reconstruction.
\end{exercise}

\begin{exercise}
Write a test that deletes the index file between two runs of the \gocode{status} logic, in process, and checks that the lines produced are identical. Then write a deliberately wrong version of \gocode{Open} that trusts a stale index, and show that the test fails.
\end{exercise}

\begin{exercise}
Run thirty copies of \gocode{golab index} at once on a ledger that is simultaneously receiving new records from \gocode{verify}. Report whether any index file ever fails \gocode{Open}, and explain why the answer does not contradict the single-ledger fork of Chapter~\ref{ch:writer}.
\end{exercise}

\begin{exercise}
Measure the time taken by \gocode{index -rebuild} for ledgers of one thousand, ten thousand, and one hundred thousand records, and plot the results. State at what size the saving from loading an existing index exceeds the cost of writing it.
\end{exercise}


\part{Distributed Systems}
\chapter{HTTP Servers and Handlers}
\label{ch:http}

\section{From a Command to a Service}

Through Part~IV the tool was something a person ran, which started, did its work, recorded the result, and exited. A service is the same logic with a different lifetime and a different set of callers. It starts once, waits, and handles requests from other programs for as long as it is left running. Two things follow from that change and both reach back to earlier parts of the book. The first is that requests arrive concurrently, so every shared resource needs the discipline of Part~III. The second is that the callers are not the author, so everything that arrives is input from outside, and the checks that were conveniences in a command become the boundary of the system.

The standard library's \gocode{net/http} package supplies the server, and its design repays attention because it is built from the same ideas as the rest of the book. The central abstraction is the interface \gocode{http.Handler}, which has a single method, \gocode{ServeHTTP}, taking a response writer and a request. A handler is anything that can answer a request. The package also defines \gocode{http.HandlerFunc}, a function type with a \gocode{ServeHTTP} method that calls the function, so that any function of the right shape can be used as a handler. That is exactly the adapter pattern of \gocode{PolicyFunc} in Chapter~\ref{ch:interfaces}, and the standard library uses it for the same reason. The service of this chapter is a single value, a \gocode{Server}, that implements \gocode{http.Handler}, and its routing table is a list of \gocode{HandlerFunc} values built from its own methods.

\golisting{internal/service/service.go}{The service: routing table, adjudication, and handlers.}{lst:service}

\section{Reading and Writing a Request}

A handler receives a \gocode{ResponseWriter} and a pointer to a \gocode{Request}, and the rules for each are strict in ways that the types do not enforce. The status line and the headers must be set before the first byte of the body is written. Calling \gocode{WriteHeader} sets the status, and writing the body without having called it sends a status of 200, so a handler that discovers an error after it has begun to write cannot change the status any more. The helper \gocode{writeJSON} in the listing therefore sets the content type, sends the status, and only then encodes the body, and it ignores the error from the encoder because nothing useful can be done once the status is committed.

The body of a request is a stream and not a value, and it comes from a client that may be slow, hostile, or simply mistaken about how much it is sending. A handler that decodes the whole body into memory without a limit lets any client exhaust the memory of the process by sending a large one. The function \gocode{http.MaxBytesReader} wraps the body in a reader that fails once a limit is passed, and the admission handler sets the limit at one mebibyte. When the limit is passed the decoder's error is an \gocode{*http.MaxBytesError}, which the handler recognizes with \gocode{errors.As}, the technique of Chapter~\ref{ch:errors}, and answers with status 413 instead of reporting a generic failure. The decoder is also told to refuse unknown fields, which is the strict choice weighed in the last exercise of Chapter~\ref{ch:encoding}. A client that misspells a field name is told so, and does not have its record accepted with the field silently dropped.

\section{One Goroutine per Request}

The server starts a goroutine for every connection and calls the handler in it. This is the model that makes Go servers simple to write, and it is the reason the concurrency of Part~III is not optional here. Two requests can be inside \gocode{admit} at the same moment, and both will want to append to the ledger. The \gocode{Server} holds a \gocode{ledger.Writer}, so both appends go through the channel to the single goroutine that owns the file. The test \gocode{TestConcurrentClientsShareOneLedgerSafely} sends forty simultaneous writes while ten other clients read, and checks afterward that the ledger holds exactly forty records and that its chain verifies. It was run under the race detector.

The readers raise a point that Part~III did not need. The handlers for the head of the ledger, for a listing of records, and for an index lookup all need the contents of the ledger. If a handler called \gocode{ledger.Load} directly, it would read the file while the owning goroutine might be in the middle of writing a line, and could see a partial record, which would fail to decode. The \gocode{Writer} therefore has a method, \gocode{Snapshot}, that sends a request of a different kind through the same channel. The owning goroutine handles it between appends, loads the file, and replies with the records. A reader sees the ledger as it stood between two appends, and never a half-written line. The test \gocode{TestSnapshotSeesCompletedAppendsAndNeverATornRecord} runs eight appenders and a repeated reader at once, verifies the chain of every snapshot, and checks that no snapshot is shorter than an earlier one. The cost is that every read loads the whole file and queues behind pending appends, which is acceptable for a teaching service and would be the first thing to change in a busy one. The protection extends only to readers inside this process. A command-line program that reads the file at the same moment has no such guarantee, which is a case of the single-process limit that Chapter~\ref{ch:writer} stated and that Chapter~\ref{ch:containers} returns to.

\section{Timeouts and Shutdown}

A server that waits forever for a client is a server that a few slow clients can stop. The \gocode{http.Server} type has timeouts for reading the request headers, for reading the whole request, for writing the response, and for idle connections, and all of them default to no limit. The command sets each of them. Without them, a client that opens connections and sends nothing holds a goroutine and a socket per connection indefinitely, and the process eventually runs out of one or the other.

Stopping a service has the same structure as stopping the pool of Chapter~\ref{ch:contexts}, with requests in place of files. The command listens for the interrupt and terminate signals through \gocode{signal.NotifyContext}, which returns a context that is cancelled when one arrives. When the context ends, the command calls \gocode{Shutdown} on the server with a deadline. The method closes the listeners so that no new connections are accepted, waits for the requests in progress to finish, and returns once they have or the deadline passes. Only then does the command close the ledger writer, and the order is essential. A writer closed first would make every request still in flight fail with \gocode{ErrClosed} on its append, after the client had been told nothing, whereas a writer closed last has every accepted append written before the process exits. The transcript in Listing~\ref{lst:session7} shows the sequence at the end, with the messages \texttt{shutting down} and \texttt{stopped} and a clean exit status.

\section{Testing Without a Network}

The package \gocode{net/http/httptest} supplies what is needed to test a handler without a real network. \gocode{NewRecorder} gives a response writer that captures what the handler wrote, and \gocode{NewServer} starts a real server on a loopback port and returns its address. The tests of this chapter use the second form, because they want the whole path, including the router, the adjudication, and the concurrency of real connections. They are collected in the next chapter, which is concerned with what they establish.

\section*{Exercises}

\begin{exercise}
Write a middleware function, one that takes an \gocode{http.Handler} and returns another, that records the status code and duration of every request, using a wrapper around the \gocode{ResponseWriter} that captures the status. Explain why the wrapper must forward \gocode{WriteHeader} and why a handler that never calls it needs a default.
\end{exercise}

\begin{exercise}
Remove the \gocode{ReadHeaderTimeout} from the server and write a program that opens many connections and sends only a few bytes on each. Report what happens to the number of goroutines in the server, and restore the timeout.
\end{exercise}

\begin{exercise}
Change the order of shutdown so that the writer is closed before the server. Start a client that posts continuously, stop the server with a signal, and report what the client observes.
\end{exercise}

\begin{exercise}
Replace the full-file read in \gocode{Snapshot} with a request that returns only the records after a given sequence number, and have the listing handler use it. State what the writer must now assume about the part of the file it does not read.
\end{exercise}

\chapter{Routing, Adjudication, and Admission}
\label{ch:adjudication}

\section{Three Questions Again}

Chapter~\ref{ch:inscription} separated three questions that a small program tends to answer in one place: whether a record is well formed, whether the actor is entitled to request the change, and whether the change is made durable. The service has a parallel set of three, and the temptation to merge them is just as strong, because a web framework makes it easy to write a single function that looks at the request, decides whether the caller may do what it asks, and does it. The service keeps them apart in its structure. \emph{Routing} chooses which handler a request is for, from the method and the path. \emph{Adjudication} decides whether the caller may use that handler. \emph{Admission} is the pipeline of Part~II, which decides what becomes of the content of a record, and it happens inside a handler, after the first two have been settled. The three have different inputs, different authorities, and different failure statuses, and the code is arranged so that each is a separate piece.

\section{Routing Decides From the Request Line}

\golisting{internal/service/route.go}{A routing table that sees only the method and the path.}{lst:route}

The router is a table of routes, each naming a method, an exact path, a handler, and the requirement a caller must meet. Its method \gocode{Match} takes two strings, the method and the path, and returns a route or a status code explaining the absence of one. The signature is the point. A function that is given only the method and the path cannot read a header, a credential, a query parameter, or the body, and so routing cannot come to depend on who is asking or what they sent. The compiler guarantees that much. That nothing else leaks into the choice, for instance through a package variable, is a matter of reading the thirty lines of the function.

When no route matches, the router distinguishes two situations that clients need to tell apart. If no route has the path, the answer is 404. If some route has the path but none with the method, the answer is 405, accompanied by an \texttt{Allow} header that lists the methods the path does support, in sorted order so that the header is the same on every request. The test \gocode{TestRouterDecidesFromMethodAndPathOnly} checks a table of such cases, including a trailing slash, which is a different path and receives 404, since the table matches paths exactly.

\section{Adjudication Is a Function of Two Values}

\golisting{internal/service/auth.go}{Bearer tokens, principals, and the permission rule.}{lst:auth}

Adjudication has two steps, which the code keeps distinct. \emph{Authentication} establishes who is calling. The service reads a bearer token from the \texttt{Authorization} header and looks it up in a list of configured tokens, each of which names a principal and says whether it may write. \emph{Authorization}, the permission rule, asks whether that principal may do what the route requires. The function \gocode{Permit} takes the principal, a flag saying whether authentication succeeded, and the route's requirement, and returns a status and a verdict. It reads nothing else. A route's requirement is one of three values: nothing, which makes the route public, as the health check is; read, which needs any authenticated principal; and write, which needs one allowed to write. An undeclared requirement is refused, for the reason that the policy vocabulary of Chapter~\ref{ch:interfaces} fails closed on an undeclared verdict: a system that decides what is permitted should deny what it does not recognize.

Two details of authentication come from the standard library and are worth stating. The comparison of tokens uses \gocode{subtle.ConstantTimeCompare}, which takes the same time whether the first byte or the last differs, so that a client cannot learn a token a byte at a time by measuring how long the server takes to refuse a guess. The loop also compares against every configured token and does not stop at the first match, so that the time does not reveal the position of the token in the list. The parsing of the configuration rejects an empty token and a token assigned to two principals, because a shared credential would make the identity of the caller ambiguous.

\begin{proposition}\label{prop:permitted}
For every request that reaches a handler of the routing table, \gocode{Permit} returned success for the principal derived from the request and the requirement of the route that matched.
\end{proposition}

\begin{proof}
Every request enters through \gocode{Server.ServeHTTP}. The only call that runs a route's handler is the last statement of that method, and it is reached only after \gocode{Match} has returned a route and \gocode{Permit} has returned success for that route's requirement. Any other path through the method returns after writing an error.
\end{proof}

The proof depends on facts about the code that the compiler does not check. The handlers are methods of \gocode{Server}, and another function in the package could call one directly and skip the check, and the value that \gocode{New} returns exposes no handler outside the package. The argument is therefore local to the package and holds as long as \gocode{ServeHTTP} remains the only way in. The tests supply evidence of the property at its most important point. A principal who may only read posts a valid batch and receives 403, an anonymous caller posts and receives 401, and in both cases the ledger afterward has no records, so the handler, which would have appended one, did not run.

\section{Routing Comes First}

The order of the steps in \gocode{ServeHTTP} is a decision with consequences. Because routing precedes adjudication, a request for a path that does not exist receives 404 whether or not it carries a valid credential, and a request with the wrong method receives 405 with an \texttt{Allow} header under every credential. The test \texttt{Test\-Routing\-Errors\-Do\-Not\-Depend\-On\-Credentials} checks both statements across four different credentials. The alternative, in which every request is authenticated before anything else, would answer unauthenticated requests with 401 even for paths that do not exist, which reveals nothing about the table. The service takes the first order because it makes the router's behaviour a function of the request line alone, and the cost is that an anonymous caller can learn which paths exist by observing which ones answer 401 or 403 and which answer 404. For a service whose routes are published this costs nothing, and for one whose routes are secret it would be a reason to reverse the order.

\section{Admission Is a Different Judgment}

\lstinputlisting[style=shellstyle,caption={A session with the service: refusal, admission, and shutdown.},label={lst:session7}]{session7.txt}

The handler for \gocode{POST /v1/admit} runs only for a caller who may write, and what it does is the work of Part~II. It decodes a batch of records, inscribes each, validates the batch as a whole, and admits the valid set under the policy the server was started with. The caller of the transcript uses the token of a principal who may write, and the server was started with the policy that denies the author \texttt{bo}. The response says that record \texttt{a} was admitted and \texttt{b}, which was written by \texttt{bo}, was refused, and gives the receipt of each. The adjudication let the caller through, and the policy still refused one of the records, which shows the separation working: the caller is entitled to submit, and the content is judged on its own terms.

A batch that fails validation, for instance because a record refers to one that is not in the batch, is answered with status 422 and is recorded in the ledger as a failed admission, as the command-line version records its own failures in Chapter~\ref{ch:inscription}. The tests cover malformed JSON and unknown fields, which are answered with 400 and are not recorded, since no batch was ever formed, and an oversized body, which is answered with 413. The distinction between a request that never became a batch and a batch that failed validation is one a reader of the ledger can rely on.

The ledger record for an admission carries the name of the principal. The first limit stated in Chapter~\ref{ch:ledger} was that the tool had no notion of identity, and here it acquires a limited one. What the ledger now records is the name attached to the token that made the request. It does not identify a person, since anyone holding the token is that principal, and the token travels as a plain header. The service as written speaks plain HTTP. A bearer token sent over a network that anyone can observe is observable, so the service belongs behind a terminator of transport encryption, or on a loopback interface, and the book does not claim more for it than that.

Refusals at the adjudication stage are not recorded. A request that receives 401 or 403 leaves no mark in the ledger. The omission is deliberate and has a cost. A ledger that recorded every refused request would let any anonymous client append to it at will, filling it with entries and making the cost of an attack the cost of the ledger's growth. The question of what the ledger should say about refusals, raised at the end of Chapter~\ref{ch:ledger}, therefore has two different answers according to who is refused. A refusal by policy, which concerns a record from an authenticated caller, is recorded as part of the admission. A refusal of the caller is not, and an operator who needs it would use the server's logs, which Part~VI takes up.

\section*{Exercises}

\begin{exercise}
Add a route for \gocode{DELETE} on the ledger and give it a requirement that no principal can meet. Show with a test that the handler is unreachable, and explain what the 403 reveals to a caller that a 404 would not.
\end{exercise}

\begin{exercise}
Reverse the order of routing and adjudication so that every request is authenticated first. Rewrite the two tests of this chapter that depend on the order, and state which property of the router is lost.
\end{exercise}

\begin{exercise}
Add a rate limit on failed authentication per remote address, recording the address and the count in memory. Discuss what the limit does to a legitimate client behind a shared address, and why the limit is state that a restart loses.
\end{exercise}

\begin{exercise}
Write a handler that records refused callers in the ledger but caps the recording at a fixed number per minute, writing a single summary record when the cap is exceeded. Decide what the summary should say and whether it can itself be used to harm the ledger.
\end{exercise}

\chapter{Containers and the Edge of the Process}
\label{ch:containers}

\section{What a Container Is}

The service runs on the machine where it was started, with that machine's files, network, and users visible to it. A container changes what the process can see and use without changing the process itself. It is not a virtual machine, which runs its own kernel, but an ordinary process of the host that the kernel has been told to treat specially. Three features of the Linux kernel do most of the work, and the others, such as the image format and the tooling, are conveniences built on them.

The first is the \emph{namespace}. A namespace gives a process its own view of one kind of global resource. A process in a new process-identifier namespace sees itself as process number 1 and cannot see the processes outside. A new host-name namespace lets it set a host name that the rest of the machine does not share. A new network namespace gives it an empty set of network interfaces, and a new mount namespace gives it a private table of mounted filesystems. The program \texttt{unshare} runs a command in new namespaces, and the transcript below was produced with it, on a host running a Linux kernel.

\lstinputlisting[style=shellstyle,caption={Namespaces demonstrated with \texttt{unshare}.},label={lst:session8}]{session8.txt}

The transcript shows each kind in turn. In a new process namespace the listing of processes contains only the command itself, as number 1, and the namespace identifiers read from the process's own entry differ between the host and the child. In a new host-name namespace the name can be changed and the host's name is untouched. In a new network namespace the only interface is the loopback. In a new mount namespace a filesystem mounted inside is not visible outside, and the final count of matching entries on the host is zero. These are the isolations that container tools combine.

The last two entries of the transcript show a consequence that matters to the service. An ordinary process that receives the terminate signal and has not arranged to handle it is ended. The first process of a process namespace is treated differently by the kernel, which delivers to it only the signals for which it has installed a handler, and a signal that cannot be caught and the kill signal are the exceptions. The same command that was terminated outside the namespace was still running after the terminate signal inside it, and was ended only by the kill signal. A service that runs as the first process of a container and does not handle the terminate signal will ignore the request to stop that container tools send first, and will be killed after a grace period without having finished anything. The \gocode{serve} command installs a handler for the signal through \gocode{signal.NotifyContext}, as Chapter~\ref{ch:http} described, and the behaviour shown here is the reason that handler is necessary and not merely tidy.

The second feature is the \emph{control group}, which limits what a group of processes may consume. A control group can cap memory, processor time, and the number of processes, and the kernel enforces the cap by reclaiming memory, by throttling, or by refusing to create more processes. A namespace limits what a process can see, and a control group limits what it can use, and the two are independent. The third feature is the root filesystem. A container starts with a directory tree of its own, assembled from the layers of an image, as its view of the files. A container whose tree contains only one program and the files that program needs has nothing else to run, to read, or to be tricked into running.

\section{A Static Binary}

The last point suggests how little a Go service needs. A Go program compiled with the \gocode{CGO\_ENABLED=0} setting contains its own runtime and its own implementation of the system calls it uses, and links to no shared library. The \texttt{file} command reports it as statically linked, and such a program runs in a directory tree that has no C library, no shell, and no certificate store. The setting has a cost for programs that depend on the C library's behaviour, such as name resolution through the system's configured services, where the pure-Go implementation takes over, and for a service that does none of these things it costs nothing. The flags \gocode{-trimpath} and \gocode{-ldflags='-s -w'} remove the build directory's path and the symbol tables from the binary, which makes it smaller and its contents less dependent on where it was built.

\lstinputlisting[style=shellstyle,caption={The image definition.},label={lst:dockerfile}]{../golab/deploy/Dockerfile}

The definition in Listing~\ref{lst:dockerfile} starts from \texttt{scratch}, which is an empty filesystem, copies in the binary, and adds an empty directory for data owned by a numeric user. The user is numeric because the image has no file of user names to look a name up in. The entry point runs the program with the ledger placed in the data directory, and the default command starts the service listening on all interfaces, since an address on loopback inside a container would be reachable only from inside it. The binary is built outside the image and copied in, by the script \texttt{deploy/build.sh}, which sets the build flags and then builds the image. A multi-stage definition that compiles inside a Go image is a common alternative and produces the same contents.

\section{Running It}

\lstinputlisting[style=shellstyle,caption={Building and running the service in a container.},label={lst:session9}]{session9.txt}

The transcript in Listing~\ref{lst:session9} follows the image from build to stop. The binary is a statically linked executable of about six megabytes, and the image built from it is under nine. The container is started with the root filesystem read-only, a memory limit of 64 megabytes, half a processor, all Linux capabilities dropped, a limit of 64 processes, and a named volume at the data directory. Each of those restricts the service further and none is needed for it to work. The service starts, logs the address it listens on, and answers an authenticated request to admit a record. The process listing from the host shows the program running as the numeric user, with a process number that belongs to the host's namespace and not the number 1 it has inside. Reading the control group's file for the memory limit on the host shows 67108864 bytes, the 64 megabytes that were asked for. That file's location is specific to the older version of control groups that this host runs, and a host on the newer version keeps the limit in a differently named file under a different path.

Stopping the container sends the terminate signal to the program, and the transcript shows the result of the handler described in the previous chapter: the stop takes a quarter of a second, the exit code recorded by the container runtime is zero, and the log ends with the messages \texttt{shutting down} and \texttt{stopped}. Had the program ignored the signal, the stop would have waited out the runtime's grace period, ten seconds by default, and ended the process with the kill signal. The last command starts a second container from the same image against the same volume, runs the audit, and reads the same ledger with one record in it. The volume is what carries the data past the life of the container, and the container itself was already stopped when the audit ran.

\section{What a Container Does Not Decide}

A container packages a program and limits what it can see and consume. It does not decide what the program is permitted to do, and the separation matters more here than anywhere else in the book, because a container looks like a security decision and is only a deployment one. The token in the environment of the service in the transcript is what lets a caller write, and the container neither knows nor checks it. The limits on memory and processes bound the damage a defect or an abusive caller can do to the host, and they do not make a request legitimate. The restrictions on capabilities and on the root filesystem reduce what an attacker who compromises the program could do next, and they do not prevent the compromise. All containers on a host share its kernel, so a flaw in the kernel is a flaw in every container's isolation, which is the difference from a virtual machine that the first section began with.

The container also leaves the single-writer limit of Chapter~\ref{ch:writer} exactly where it was. The service is the sole owner of the ledger within its process, and it is the sole owner of the file only if nothing else appends to it. Two containers started from the image against the same volume would each have a writer on the same file, and they would fork the chain just as twenty command-line processes did. The audit container of the transcript only reads, and reads were the one thing a second process could do safely, up to the caveat in Chapter~\ref{ch:http} that it might see a line mid-write. A deployment of this service is therefore a single instance, and scaling it by running more copies is not available without the file lock of the exercise in Chapter~\ref{ch:writer} or a different arrangement for the ledger altogether.

\section{Where Part V Ends}

The tool is now a service with a defined interface, a routing table that is a function of the request line, a permission rule that is a function of the caller and the route, an admission step that is the pipeline of Part~II, a ledger owned by one goroutine, and a container that runs it with limits. Each layer has been given a statement of what it guarantees and what it does not, and Appendix~\ref{app:guarantees} collects them. Part~VI turns from building the tool to measuring it. It asks how fast the ledger and the pool are, what the service is doing while it runs, how the tool behaves when its disk or its network fails, and whether a compressed summary of a ledger can still support the decisions the full ledger could.

\section*{Exercises}

\begin{exercise}
Run the service in a container that has no handler for the terminate signal, by removing the \gocode{signal.NotifyContext} call, and time the stop. Report the time and the exit code, and explain both in terms of the first process of a namespace.
\end{exercise}

\begin{exercise}
Start the container with the root filesystem writable and the data directory not mounted as a volume. Add a record, stop and remove the container, start another, and report what the ledger contains. State what this shows about where state belongs.
\end{exercise}

\begin{exercise}
Run two containers against the same volume, each posting to its own service, and report what the audit says afterward. Then propose an arrangement, short of a file lock, in which only one of them ever writes.
\end{exercise}

\begin{exercise}
Write a multi-stage image definition that compiles the program inside a Go build image and copies the result into the empty image. Compare the final image with the one built here, and say what the extra stage adds to the time of a build.
\end{exercise}


\part{Experimental Systems}
\chapter{Measurement}
\label{ch:measurement}

\section{Performance Claims Are Claims}

Earlier chapters made statements about cost without measuring anything. Chapter~\ref{ch:writer} said that each append to the ledger reads and checks the whole file, so that the cost of $n$ appends grows as $n^2$, and Chapter~\ref{ch:index} said that an index saves work when it is extended more often than it is rebuilt. Both statements follow from reading the code, and both are the kind of statement that a measurement can confirm, refine, or contradict. Part~VI treats the tool as an object of study, and the first instrument is the benchmark. A statement about speed deserves the same treatment as any other claim in the book: it should say what was measured, on what, with what result, and what the result does not show.

The \gocode{testing} package supports benchmarks as ordinary functions. A function whose name begins with \gocode{Benchmark} and which takes a \gocode{*testing.B} is run by \gocode{go test -bench}, and its body loops \gocode{b.N} times. The framework chooses \gocode{b.N}: it runs the function with a small value, measures, and raises the value until the total time is long enough to be reliable, then reports the time per iteration. Inside the function, \gocode{b.StopTimer} and \gocode{b.StartTimer} exclude setup work from the measurement, and \gocode{b.Run} creates sub-benchmarks, which is how one function measures a range of sizes. The flag \gocode{-benchmem} adds the number of bytes and allocations per iteration, and \gocode{b.SetBytes} lets the framework report throughput.

\section{The Cost of an Append}

\golisting{internal/ledger/bench_test.go}{Benchmarks for appending to a ledger.}{lst:bench-ledger}

The benchmark of Listing~\ref{lst:bench-ledger} measures one append to a ledger that already holds $n$ records, for $n$ equal to a hundred, a thousand, and ten thousand. To keep $n$ constant, it restores the file to its starting contents before every iteration, with the timer stopped, so that each measured append sees a ledger of exactly the stated size. The results come from one run on a machine with two processors.

\lstinputlisting[style=shellstyle,caption={Benchmarks of the ledger, the index, and the pool.},label={lst:session11}]{session11.txt}

The three results for appending are 1.35 milliseconds, 10.6 milliseconds, and 102 milliseconds. Each tenfold increase in the size of the ledger multiplies the time by roughly ten, and the number of allocations grows in the same proportion, from about two thousand to about two hundred and forty thousand. The cost of one append is proportional to the number of records already present, which is the behaviour Chapter~\ref{ch:writer} predicted. The cost per record is about ten microseconds on this machine. If the trend continued, a ledger of a million records would take about ten seconds to append to once, and a ledger of a million records built by a million appends would have cost the sum, which is quadratic. The extrapolation is a statement about a straight line drawn through three points. It is not a measurement of a ledger of a million records, and the measurements here establish only that the growth is linear over the range from a hundred to ten thousand.

A benchmark says how much and not why. The profiler says where.

\section{Where the Time Goes}

The standard toolchain can sample the running program and report which functions the samples fall in. The flag \gocode{-cpuprofile} on \gocode{go test} writes a profile, and \gocode{go tool pprof} reads it. The report distinguishes \emph{flat} time, spent in a function's own instructions, from \emph{cumulative} time, spent in the function and everything it calls, and for finding the structure of a cost the cumulative column is the useful one.

\lstinputlisting[style=shellstyle,caption={Selected rows of a CPU profile of appending to a ledger of ten thousand records.},label={lst:session12}]{session12.txt}

The profile in Listing~\ref{lst:session12} says that \gocode{Append} accounts for about four fifths of all the samples, the remainder being the restoration of the file between iterations. Within \gocode{Append}, loading the ledger takes the larger share, 55 percent of all samples, and nearly all of it is spent decoding JSON. Checking the chain takes 27 percent, and almost all of that is recomputing digests. The write of the new record does not appear among the visible costs. The time is going into reading and verifying the ledger that is already there, which is the work that the single-writer design of Chapter~\ref{ch:writer} was said to repeat. A writer that kept the head digest and the next sequence number in memory, and checked the file only when it opened it, would remove both of the large items from the per-append cost. That change is the third exercise of Chapter~\ref{ch:writer}, and it has a price, since the writer then trusts its own memory between appends. The measurements identify what could be saved and say nothing yet about whether the saving is worth the price.

\section{The Flush to Stable Storage}

Part~VI adds one thing to \gocode{Append}, a call to \gocode{Sync} on the file before it is closed, so that the record is on stable storage when the function returns and not merely in the operating system's cache. The second benchmark compares an append with and without the call, on a ledger of a hundred records, where the fixed costs dominate. The result is 1.74 milliseconds with the flush and 1.00 millisecond without, so the flush adds about 0.74 milliseconds to each append on this machine. That number belongs to this machine's storage and not to the program. A different device would give a different figure, and a filesystem in memory would give almost none. The decision to flush is a decision about durability: without it, a crash can lose records that were reported as written. For a ledger whose purpose is to be relied on after the fact, the book takes the cost, and the benchmark makes the size of that cost visible, which a decision of this kind should have.

\section{The Index and the Pool}

The index benchmarks measure two things. Building an index from a whole ledger in memory takes 3.1 milliseconds for a thousand records, 34.8 for ten thousand, and 380 for a hundred thousand, which is between three and four microseconds per record and again linear. Applying a single record to an existing index takes about 118 nanoseconds and allocates nothing. The ratio between the two figures is large, and it needs care. The cost of 118 nanoseconds is the cost of \gocode{Apply} alone. The method \gocode{CatchUp}, which is what \gocode{Open} calls, checks the whole chain before applying new records, as Chapter~\ref{ch:index} noted, and a check of the chain is most of what a build costs. In the present code, loading and extending an index is therefore not as cheap as the figure for \gocode{Apply} suggests, and the fourth exercise of Chapter~\ref{ch:index}, which proposes checking only the new records, is where that saving would come from.

The pool benchmark hashes sixty-four files of one mebibyte with one, two, four, and eight workers. The times fall from 255 milliseconds to 144, 133, and 117, and the throughput rises from 263 to 465, 506, and 576 megabytes per second. The machine has two processors, and the step from one worker to two gives a speedup of 1.8, which is what a CPU-bound task on two processors should give. The further gains from four and eight workers are smaller, and the figure for eight is above what two processors could deliver for pure computation. The measurement does not say why. It may reflect time the workers spend waiting on the operating system, which other workers can overlap, or noise in a single run on a shared machine. A claim about the cause would need further experiments, and the measurement supports only the statement that more workers than processors still helped a little in this run.

\section{What a Benchmark Cannot Say}

Every number in this chapter comes from one run on one virtual machine whose other tenants are unknown. The benchmark framework reports a mean and no spread. A responsible comparison repeats each benchmark several times with the \gocode{-count} flag and compares the distributions, with a tool such as \texttt{benchstat} from the Go project, and a single figure cannot show whether a difference between two figures is larger than the run-to-run variation. The files in the pool benchmark fit in the operating system's cache after the first pass, so the figures describe hashing and scheduling and say nothing about a cold disk. The environment that produced the numbers is printed at the head of the benchmark output, with the operating system, the architecture, and the processor, and it belongs next to any figure that is quoted. A number without its environment is not a measurement of anything in particular.

\section*{Exercises}

\begin{exercise}
Implement the writer that holds the head digest and the next sequence number in memory, and repeat the benchmark of this chapter for ledgers of a hundred, a thousand, and ten thousand records. Report the new figures, the new profile, and the new risks that the cache introduces.
\end{exercise}

\begin{exercise}
Run the benchmark of the pool ten times with the \gocode{-count} flag and report the range of each figure. State which of the conclusions of this chapter survive the variation and which do not.
\end{exercise}

\begin{exercise}
Write a benchmark of \gocode{ledger.Load} alone for the three sizes, and compare its time with the share the profile assigned to loading. Explain any discrepancy.
\end{exercise}

\begin{exercise}
Replace the call to \gocode{Sync} with a policy that flushes once every hundred appends or every second, whichever comes first. Describe exactly what a crash can now lose, and decide whether a ledger with such a policy can still make the statement that a record reported as written is durable.
\end{exercise}

\chapter{Logs and Faults}
\label{ch:faults}

\section{Two Records With Different Jobs}

The service now keeps two kinds of record, and they should not be confused. The ledger is the durable account of what the tool decided and did. It is chained, it is checked, and it refuses to grow when damaged. A log is an operational stream, written for the person who runs the service so that they can see what is happening and diagnose what went wrong. Nothing in the log is chained, lines can be lost when a process dies, and the log is not evidence of anything in the sense the ledger tries to be. The distinction decides what goes where. A refusal of an anonymous caller does not belong in the ledger, for the reason given in Chapter~\ref{ch:adjudication}, that anyone could fill the ledger by being refused, and it does belong in the log, which can be rotated and discarded.

The standard library's \gocode{log/slog} package writes structured entries, each with a level, a message, and a list of named values, and its handlers render them as JSON lines or as readable text. The service holds a \gocode{*slog.Logger}, and \gocode{ServeHTTP} writes one entry per request after the handler has finished, whatever the outcome, with the method, the path, the status, the number of bytes sent, the duration in milliseconds, and the name of the principal. To know the status after the fact, the method wraps the response writer in a small recorder that remembers the status that was sent, and the recorder implements the \gocode{Unwrap} method so that code reaching for optional features of the underlying writer, such as flushing, can still find it. The level follows the status: server errors are logged as errors, client errors as warnings, and everything else as information.

\lstinputlisting[style=shellstyle,caption={Three requests and the entries they produced.},label={lst:session13}]{session13.txt}

The entries of Listing~\ref{lst:session13} show the three cases. A reader may fetch the head of the ledger, a reader who tries to write is refused with status 403, and a caller with an unknown token is refused with 401, with an empty principal because authentication failed. The refusals appear at the warning level, which is where an operator looks for them, and nothing about them reaches the ledger. The entries contain no credential. That is a property worth testing and not merely intending, because a log is the place where secrets leak: the test sends requests with a valid token and an invalid one and checks that neither string appears in the output, nor the word naming the authentication scheme. The query string is also left out of the entry, since a caller might put anything in it. What is logged is chosen, field by field, and the default is to leave a value out.

The same discipline applies to what the service tells the client. The handlers that fail because of the server's own machinery, such as a ledger that cannot be written, now answer with the status 500 and the words \texttt{internal error}, and send the cause to the log. A message from the operating system names file paths, and a client has no use for the layout of the server's disk. The test of this behaviour points the writer at a directory that does not exist, posts a batch, and checks that the response contains no part of the path while the log entry does.

\section{Faults as Inputs}

A program's behaviour on success is the part its author thinks about. Its behaviour when the machine fails to cooperate is mostly unexamined until it happens, and the failures that matter most to a ledger, a full disk, a write that does not complete, a flush that fails, are the ones that tests do not produce on their own. Fault injection is the practice of producing them on purpose. There are two ways to do it, and the chapter uses both. A simulated fault replaces part of the environment with something that fails in a chosen way, and it is repeatable and fast. A real fault arranges for the environment itself to fail, and it is more convincing and more awkward.

The real fault used here is a full filesystem. In a private mount namespace, as in Chapter~\ref{ch:containers}, the experiment mounts a small memory-backed filesystem of eight kibibytes and appends records to a ledger on it until an append fails. The ledger as it stood at the end of Part~V behaved as follows.

\lstinputlisting[style=shellstyle,caption={Appending until the filesystem is full, before and after the repair.},label={lst:session10}]{session10.txt}

In the first half of Listing~\ref{lst:session10} the thirtieth append failed with the error that the device had no space, which is the correct report. But the file was found to be exactly 8192 bytes, the whole of the filesystem, and the audit failed with an error about line 30, ending unexpectedly. The write of the thirtieth record had succeeded in part. The first bytes of the line fitted in the last block of the file, the rest did not, and the operating system returned an error after storing the fragment. The ledger now ended in half a record. Every later command that loaded the ledger failed, so the failure of one append had destroyed the usability of the whole ledger. The protection built into \gocode{Append}, which refuses to extend a broken chain, did what it was designed to do, and what it protected was a ledger that the failed append had itself broken.

The repair, in the second half of the listing, is to remember the length of the file before writing and to cut the file back to that length if the write does not complete. The same experiment now ends with the same twenty-nine successful appends and the same error on the thirtieth, a file of 8134 bytes, and an audit that finds twenty-nine records and a valid chain. The failed append left no trace in the file. The fix is a few lines in the function that writes the line.

\section{Seams and Simulated Faults}

The real fault cannot be part of the ordinary test suite, since it needs privileges and a mount, so the repair has tests of the simulated kind. The ledger package declares a small interface, \gocode{appendFile}, naming the five methods of an open file that the writing of a record uses, and a package-level variable, \gocode{openAppend}, that returns a value of that type. In normal use the variable opens a real file. A test replaces it with a function that returns a wrapper around a real file, which can be told to write only part of a line and then fail, to refuse to truncate, or to fail when synchronizing.

\golisting{internal/ledger/fault_test.go}{Simulated faults for the ledger.}{lst:fault-test}

A seam of this kind costs something, and the cost should be stated. A package-level variable that tests replace is shared state, so two tests that replace it cannot run in parallel, and a test that forgets to restore it affects the tests that follow. The interface is wider than the code strictly needs, and it exists only for the tests. The alternative is to pass the opener as a parameter to every function, which changes every caller for the sake of a test. Each choice trades something, and the book takes the variable because the ledger has one write path and the number of tests that touch it is small. A larger system would likely choose differently.

The first test makes the wrapper write forty bytes of a record and fail, then checks three things: that \gocode{Append} returned the disk error, that the file's bytes are identical to what they were before the call, and that a later append succeeds and the ledger loads with a valid chain. The second makes the truncation fail as well, and checks that the error reports both failures, saying that the partial record could not be removed, and that the damaged ledger is visible to \gocode{Load}, which is the honest outcome when the damage cannot be undone. The third test is the most important for the contract of the function.

\section{An Error Is Not a Statement That Nothing Happened}

The third test makes the synchronization fail after a complete write. \gocode{Append} returns an error, since it cannot promise that the record is on stable storage, and the record is in the file. The test checks that the ledger contains it. A function that returns an error has told its caller that something went wrong, and the usual reading is that the operation did not take effect. That reading is false here, and it is false for a large class of operations that end with a step that can fail after the effect has occurred. A caller that retries after such an error will write the record twice.

The documentation of \gocode{Append} says so in a sentence: an error means that the record is not guaranteed to be present, and not that it is absent. The service passes the consequence on to its clients in the only way it can, by answering a failed append with a failure status and recording the cause in the log. A client that receives that status does not know whether its batch was recorded. The design that would let a client know is one in which the client supplies an identifier for the request and the server remembers which identifiers it has completed, so that a retry can be recognized, and the fourth exercise asks for it.

The repair also has a limit. If the process or the machine stops between the partial write and the truncation, the fragment remains, and the next command to load the ledger will report the broken line and stop. This is the loud failure and not the silent one, since the ledger refuses to proceed and does not pretend. Repairing such a ledger by cutting the torn final line is a decision about trust, because the same cut could hide the deletion of a real record, and the third exercise asks for a command that makes the cut only when everything before it verifies.

\section*{Exercises}

\begin{exercise}
Write a test that stops a goroutine in the middle of \gocode{writeLine}, between the write and the truncation, by using the wrapper to block, and then loads the ledger from another goroutine. Report what the loader sees and explain why the single-writer design of Chapter~\ref{ch:writer} does not prevent it.
\end{exercise}

\begin{exercise}
Add a fault to the wrapper that makes \gocode{Close} fail after a successful write and flush. Decide what \gocode{Append} should report, and whether the record can be said to be absent, present, or unknown.
\end{exercise}

\begin{exercise}
Write a command \gocode{repair} that removes a final line of the ledger that does not decode, only if the records before it form a valid chain, and records the repair in the ledger. Discuss what a repair that hides the removal of a valid record would look like, and what the command must report to avoid it.
\end{exercise}

\begin{exercise}
Add an optional request identifier to the admission endpoint. Make the server remember, for a bounded time, the identifiers it has completed and the response it gave, and answer a repeated identifier with the stored response. State what the server must do about the identifier of a request whose append returned an error.
\end{exercise}

\chapter{Manifests and Reproducibility}
\label{ch:manifests}

\section{What a Result Depends On}

A figure from an experiment is a statement about a particular run, and a reader who wants to rely on it needs to know what the run consisted of. At least three things determine a result: the inputs it was computed from, the settings that were chosen, and the program that did the computing. A fourth, the machine and its state, matters for some results and not for others. The benchmarks of the previous chapter depend on all four, which is why their output opens with a description of the machine. An experiment whose result does not depend on timing, such as the one of the next chapter, depends only on the first three, and the first three can be written down.

A \emph{manifest} is that written record. The manifest of \golab{} lists the experiment's name and its parameters, the version of the Go toolchain that built the program, the module path, the build settings that the toolchain embedded in the binary, and for each input file its path, its size, and its SHA-256 digest. The package that builds it follows two rules that make it useful as evidence. It contains no clock reading and no machine name, so that two runs under the same conditions produce byte-identical manifests, and the paths of inputs are recorded relative to the input directory, so that moving the files does not change the manifest. The time at which a manifest was made is recorded by the ledger entry that refers to it, which carries the manifest's digest. The manifest is then a function of what it describes, and the ledger supplies the fact that the description was made at a certain point in the sequence.

\golisting{internal/manifest/manifest.go}{The manifest of an experiment.}{lst:manifest}

\section{Build Information}

The toolchain embeds facts about a build in the binary, and the function \texttt{debug.\allowbreak ReadBuildInfo} reads them back. The record includes the version of the toolchain, the path of the main module, and a list of key and value pairs for the settings of the build: the compiler, the target operating system and architecture, whether the C compiler was enabled, the level of the target instruction set, and, when the source is in a version-control repository, the revision and whether the working tree had uncommitted changes. The manifest copies these into its own fields.

\lstinputlisting[style=shellstyle,caption={Writing a manifest, checking it, and checking it after the inputs change.},label={lst:session15}]{session15.txt}

The transcript in Listing~\ref{lst:session15} writes the manifest of a directory of two files, prints it, and checks it. The check recomputes the manifest from the present state and prints the differences, and an unchanged state gives the message that the manifest matches. After one file is overwritten and another added, the check names exactly those two inputs and exits with a failure status. The ledger gained one entry, the action \texttt{manifest} with the digest, and the check wrote nothing.

Two entries of the printed manifest are worth a closer look. The first is that the build was made with the C compiler enabled. The second is a long setting called \texttt{DefaultGODEBUG} that lists compatibility switches. The list is there because the module file declares an older version of the language than the toolchain implements, and the toolchain then makes some of its newer behaviours revert to the old ones. One entry of the list, \texttt{httpmuxgo121=1}, makes the standard library's request multiplexer ignore the pattern syntax that later versions added to it. The service does not use that multiplexer, since its routing table is the one of Chapter~\ref{ch:adjudication}, so the entry changes nothing for the code of this book. A service built on the standard multiplexer under the same module file would route differently from what the version of its toolchain suggests, and the manifest is where that would be found.

\section{Comparing Two Builds}

\lstinputlisting[style=shellstyle,caption={The static binary of Chapter~\ref{ch:containers} checked against the manifest of the ordinary build.},label={lst:session16}]{session16.txt}

The binary of the container image was built with the C compiler disabled and with the option that removes file-system paths from the executable. Checked against the manifest made by the ordinary build, it reports the differences in Listing~\ref{lst:session16}: the added option, the changed value for the C compiler, and the removal of the settings for the C toolchain's flags, which a build without the C compiler does not record. The two builds were meant to differ and the manifest says how.

The same listing shows a limit. The container build also passed an option to the linker that strips symbol tables, and that option does not appear in the differences. In this toolchain version the settings embedded in the binary do not include it, so a binary built with it and one built without it have equal manifests. A manifest identifies a build only up to what the toolchain chooses to record, and the missing option is a case where two different binaries are indistinguishable by their manifests. The first exercise closes that gap with the digest of the executable itself.

A second limit concerns the source. When the source is not under version control, as it is not in the directory where these transcripts were produced, the settings contain no revision, and two different source trees compiled with the same toolchain and flags yield the same manifest. The record of the revision is only as good as the habit of building from a clean checkout, and when the tree is not clean, the setting that says so is the one a careful reader looks for.

\section{What the Manifest Does Not Establish}

A manifest says what an experiment started from. It does not say that the experiment, run again from the same start, would give the same result. That is a property of the program, and it holds or fails for reasons the manifest cannot see: use of the clock, the order in which goroutines finish, the iteration order of a map, and the contents of the environment all affect the outcome of programs, and a manifest records none of them. The experiment of the next chapter has been arranged so that the property can be checked, since it is built from nothing that varies, and its tests include running the search twice and comparing. A benchmark does not have the property, and its manifest would fix the inputs and leave the numbers free to differ by the noise of the machine.

The manifest also does not establish that the inputs are the right ones. It records which files were used, by digest, and a digest is the identity of content, as Chapter~\ref{ch:traversal} said, and not a statement that the content is correct. An experiment run on the wrong files has a perfectly good manifest.

\section*{Exercises}

\begin{exercise}
Add the SHA-256 digest of the running executable to the manifest, found through \gocode{os.Executable}. Show that two builds that differ only in the stripping option now have different manifests, and say what the change does to the property that moving the program does not change the manifest.
\end{exercise}

\begin{exercise}
Build the program inside a version-control checkout and again after modifying a file without committing it. Report the settings that differ and explain how a reader of two manifests could tell the builds apart.
\end{exercise}

\begin{exercise}
Make the check command ignore the settings and compare only the inputs and parameters. Describe a result that this would let through that the full check catches.
\end{exercise}

\begin{exercise}
Write the manifest of the benchmark run of Chapter~\ref{ch:measurement}, including the machine description from the benchmark header as parameters. Run the benchmark twice and compare the manifests and the results, and state what the manifests establish about the difference between the results.
\end{exercise}

\chapter{What a Summary Discards}
\label{ch:compression}

\section{The Question}

A ledger grows without bound, and sooner or later someone proposes to replace the old part of it with something smaller. The index of Part~IV is one such replacement, and it is safe because the ledger stays where it is and the index can be thrown away. The proposal that follows naturally is to keep the summary and discard the records, and the question this chapter asks is what that would cost. The question is put as an experiment on records and ledgers and has nothing to do with any machine-learning model. It concerns the plainest form of the principle that has run through the book, that a record of what happened and a judgment about it are different things, and that a summary which keeps the judgments someone expected to need may have dropped the distinctions that a later judgment requires.

\section{Summaries, Decisions, and Availability}

Fix a set of ledgers. A \emph{summary} is a function that maps each ledger to a smaller description of it. A \emph{decision} is a question about a ledger whose answer someone may later need: what was the last digest recorded for a file, did a failure follow a success, how many records are there. A decision is \emph{available} from a summary if its answer can be computed from the summary alone.

\begin{proposition}\label{prop:available}
Let $S$ be a summary and $D$ a decision on a set $\mathcal{L}$ of ledgers. $D$ is available from $S$ on $\mathcal{L}$, meaning that some function $f$ satisfies $f(S(L)) = D(L)$ for every $L \in \mathcal{L}$, if and only if no two ledgers $L_1, L_2 \in \mathcal{L}$ have $S(L_1) = S(L_2)$ and $D(L_1) \neq D(L_2)$.
\end{proposition}

\begin{proof}
If $f$ exists and $S(L_1) = S(L_2)$, then $D(L_1) = f(S(L_1)) = f(S(L_2)) = D(L_2)$. Conversely, if no such pair exists, define $f(s)$ for a summary value $s$ in the image of $S$ to be $D(L)$ for any $L$ with $S(L) = s$. The choice does not matter, since all such ledgers agree on $D$, and so $f$ is well defined and satisfies the equation.
\end{proof}

The proposition turns a question about computation into a question about pairs, and pairs can be searched for. A pair of ledgers with equal summaries and different answers is a \emph{counterexample}, and a single counterexample settles the matter for every set of ledgers that contains both. The converse is weaker. Finding no counterexample among the ledgers examined shows that the decision is available on those ledgers, and it shows nothing about longer ones unless an argument covers them.

A second fact orders the summaries. Call $S'$ a refinement of $S$ if $S'(L_1) = S'(L_2)$ always implies $S(L_1) = S(L_2)$, which means that $S'$ keeps at least what $S$ keeps. Every decision available from $S$ is then available from $S'$, because a counterexample for $S'$ would be one for $S$. A summary that adds to another never loses a decision, and the table below shows the relation in the rows for the index and the counts.

\section{The Search}

The package \gocode{summary} makes the search concrete. It fixes a small alphabet of seven events: two successful verifications of one file with different digests, one successful verification of a second file, a failed verification of each file, and a run of a program that succeeds or fails. It turns every sequence of events of length zero through six into a real ledger chain in memory, using the function \gocode{ledger.Extend}, and so the records are the genuine article, with sequence numbers, predecessor digests, and their own digests. For each ledger it computes six summaries and six decisions, groups the ledgers by summary value, and looks, for each summary and decision, for a group in which the answers differ. Lengths are visited in increasing order, so the counterexample reported for a cell is among the shortest that exist.

\golisting{internal/summary/summary.go}{The search for what summaries lose.}{lst:summary}

The six summaries are the full record, represented by the digests of all its records; the counts of records by action; the per-path index of Chapter~\ref{ch:index}, without the sequence numbers and without the head digest; the index together with the counts; the last three records; and the index together with two flags that were added after seeing the decisions. The six decisions are the last digest recorded for the first file, whether the second file is currently verified, the number of failures recorded for the first file, whether the first file failed after it had succeeded, whether a run failed since the second file's last verification attempt, and the number of records. The summaries and decisions are chosen by the author and they are a sample, not a census of the questions a reader might ask. The experiment is a demonstration that loss occurs and of what form it takes.

\lstinputlisting[style=shellstyle,caption={The search over all ledgers of up to six records.},label={lst:session14}]{session14.txt}

\section{Reading the Table}

The first row of Listing~\ref{lst:session14} is the control. The full record distinguishes all 137,257 ledgers, and every decision is available from it, as it must be. The search found no flaw in the machinery, and a test checks the same thing.

The row for the counts reduces the 137,257 ledgers to 84 distinct values and keeps one decision, the number of records. It cannot say anything about a particular file, because it counts actions and not subjects. The row for the per-path index reduces the ledgers to 611 values and keeps the first three decisions. These are the questions the index was designed to answer, and for them the experiment confirms what the structure of the index already argued: the answers are fields of the entries, so no two ledgers with the same entries can disagree about them. It loses the other three, and the counterexamples say why.

The fourth decision, whether the file failed after it had succeeded, is lost because the entry records how many attempts and failures there were and whether the last attempt succeeded, and not their order. The pair printed for it has three events in each ledger: a success, a failure, and a second success of the same digest in one, and a failure followed by two identical successes in the other. Both leave the file verified with the same digest, two successes, and one failure, and the index cannot tell them apart. In the first the failure came after a success, and in the second it did not. The fifth decision is lost for a simpler reason. The index takes no notice of runs, so a ledger that holds a failed run and a ledger that holds nothing have the same index. The sixth is lost because the index counts only verification records. The shortest counterexample for it has an empty ledger on one side and a ledger of one successful run on the other.

The row for the index with the counts keeps everything the two parts keep, as the refinement relation requires, and adds nothing that neither had. The row for the last three records shows that recency is not sufficiency. A window of the most recent events keeps no decision at all, since a ledger can have any history before the window. The last row is the instructive one. With two flags added that answer the fourth and fifth decisions, the summary makes both of them available, and the sixth decision, which it did not anticipate and for which it has no count, is lost. The two flags could have been chosen only by knowing the questions in advance. A summary can always be extended to keep the answer to a question that is known, and it cannot be extended to keep the answer to one that has not been asked.

\section{The Digest That Answers Nothing}

One summary has been left out of the table on purpose. Suppose the summary is the number of records together with the digest of the last record. Since the digest depends on every record before it, two different ledgers have different digests unless a collision exists, and the summary is therefore injective, so by Proposition~\ref{prop:available} every decision is available from it. This is correct and useless. The function $f$ whose existence the proposition promises would have to recover the ledger from its digest, and nobody can compute that function. The proposition says that an answer exists, and does not say that it is computable by any procedure that is cheaper than keeping the ledger.

The observation explains why the head digest in the index is not there to answer questions. It binds the index to one ledger, so that a mismatch is noticed, and it does so without making the index any more informative. It also marks the line between two uses of a summary. A summary that is kept alongside the ledger can be checked against it, and recomputed from it whenever there is doubt, and the doubt is cheap to resolve. A summary that is kept in place of the ledger can be checked against nothing, and the distinctions it dropped cannot be recovered when a later decision needs them.

\section{What the Experiment Does Not Show}

The experiment shows the form of the loss and does not measure its importance. A counterexample of three events establishes that the index cannot answer one question, and says nothing about whether anyone will ever ask it. The alphabet has seven events and two files, and a real ledger has many more, so the table is a small model of the situation. The cells marked \emph{no} are exact: each has a pair that anyone can check, and the test suite checks every pair afresh. The cells marked \emph{yes} for the index and the counts are backed by the argument that the answers are fields of the summary, and for the others by the search over ledgers up to six records, which is evidence for those lengths and not a proof beyond them. The search takes about fifteen seconds on the machine that produced the transcript, and its cost grows by a factor of seven with each added record, which limits how far the same method can go. The experiment is deterministic, so the manifest of Chapter~\ref{ch:manifests} that records its one parameter is sufficient to reproduce the table.

The practical lesson is the one the ordering of the chapters has prepared. Derived data is safe to throw away when the source stays. A summary that replaces its source is a bet on which questions will matter, and the experiment shows exactly where such a bet can be lost.

\section*{Exercises}

\begin{exercise}
Add a seventh summary that keeps the index and the last event, and a seventh decision that asks what the final event of the ledger was. Predict the new row and column before running the search, and compare.
\end{exercise}

\begin{exercise}
Prove that the refinement relation of this chapter is a preorder, and show that the summary of the index with the counts refines both of its parts. Use the proof to explain why the table cannot contain a decision that is available from the index but not from the index with the counts.
\end{exercise}

\begin{exercise}
Extend the alphabet with a third file and report how the number of ledgers and the running time change. State what the larger alphabet could change about the cells marked \emph{yes}, and what it could not change about the cells marked \emph{no}.
\end{exercise}

\begin{exercise}
Design a summary of a ledger that is compact, makes the fourth and fifth decisions available, and is built without knowing those two decisions in advance. Say what you gave up to do it, and test your design with the search.
\end{exercise}


\chapter*{Conclusion}
\addcontentsline{toc}{chapter}{Conclusion}

The book began with a counter that could not go below zero and ended with a search over every small ledger for the questions a summary cannot answer. Between them it built one program in six stages, and the same distinction was drawn at each stage: what the compiler guarantees, what a library guarantees under stated assumptions, and what the program merely arranges and checks by tests. The distinction is the book's main claim about the language. Go gives a programmer types that cannot be forged outside their package, interfaces that name only what a consumer needs, errors that are values to be examined, and goroutines that are cheap enough to use freely, and it leaves unenforced almost everything else: that errors are examined, that a map is not shared, that a function is pure, that a callback does not block. Good Go is mostly a matter of arranging the code so that the unenforced things become visible and testable, and of saying in the documentation which ones remain a matter of discipline.

The same distinction governs the analogies to the author's other work. The staging of inscription, validation, admission, projection, and activation, the separation of routing from adjudication, and the observation that a summary can discard what a later gate requires are structural analogies. Each was implemented in a small program and tested, and none of the formal properties of the specifications they resemble were claimed or proved here. Where a proposition appears in the text, it is a proposition about this code under the stated assumptions, and where the assumptions are not enforced, the text says so. Appendix~\ref{app:guarantees} collects the statements in one table, and its second column is the honest summary of the book: a short list of things the toolchain does for the programmer, and a longer list of things the programmer must keep doing.

Several limits remain open and are named in the text where they arise: the ledger is tamper-evident only relative to a head digest held elsewhere, two processes appending to one file can still fork it, the service speaks plain HTTP, the index file is trusted beyond its position in the ledger, and the manifest identifies a build only as far as the toolchain records it. Each could be closed with more machinery, and each closure would raise a new question of the same shape, which is a fair description of how systems of this kind are built.


\appendix
\chapter{Complete Source of the Command}
\label{app:main}

The command-line entry point is reproduced in full. Everything else in the \golab{} module appears in the body of the chapters.

\golisting{cmd/golab/main.go}{The \gocode{golab} command.}{lst:main}

\chapter{Building and Running}
\label{app:build}

The module requires a Go toolchain and no other dependency. From the directory that contains \texttt{go.mod}, the command \gocode{go test ./...} builds every package and runs every test, the command \gocode{go vet ./...} applies the standard static checks, and the command \gocode{go build -o golab ./cmd/golab} produces the executable used in the transcript of Chapter~\ref{ch:ledger}. The programs of Chapter~\ref{ch:transition} live under \texttt{examples} and are each run with \gocode{go run} followed by the path of the example directory. The service is started with \gocode{golab serve}, which reads its tokens from the environment variable \texttt{GOLAB\_TOKENS}. The script \texttt{deploy/build.sh} builds the static binary and, where a container runtime is available, the image of Chapter~\ref{ch:containers}. The commands \gocode{manifest} and \gocode{experiment compress} of Part~VI need no setup; the second takes about fifteen seconds at its default size. The benchmarks of Chapter~\ref{ch:measurement} run with \gocode{go test -run xxx -bench .} from the module root. The tests that exercise subprocesses assume that a POSIX shell named \texttt{sh} is on the path, which is the case on Linux and under the Windows Subsystem for Linux.

\chapter{An Inventory of Guarantees}
\label{app:guarantees}

The chapters of Parts~I to VI repeatedly distinguish what the Go language or toolchain guarantees from what a library guarantees under assumptions and from what the program merely arranges by convention and checks by tests. Table~\ref{tab:guarantees} collects the distinctions in one place. The table is deliberately conservative. A property is placed in the first column only if a mechanism outside the author's discipline enforces it, and a property that depends on the callers behaving correctly is listed as a convention however carefully it is documented.

\begin{table}[ht]
\centering
\small
\begin{tabular}{@{}p{0.31\textwidth}p{0.17\textwidth}p{0.44\textwidth}@{}}
\toprule
Property & Status & Qualification \\
\midrule
Code outside \gocode{artifact} cannot build a \gocode{Verified} with an invented digest & Compiler & Visibility of unexported fields. The zero value remains available to every caller. \\
A zero \gocode{Verified} is not mistaken for evidence & Convention & Callers must test \gocode{IsZero}; nothing forces them to. \\
External modules cannot import \gocode{internal} packages & Toolchain & Enforced by the \gocode{go} command. \\
Import cycles do not exist & Compiler & Checked on every build. \\
A returned error is examined by the caller & Convention & The compiler accepts a discarded error; a static analyser can flag it. \\
\gocode{Kind} takes only its four declared values & Convention & A conversion from any string compiles. \\
A policy returns only the five verdicts, never \emph{malformed} or \emph{duplicate} & Compiler & The return type has no such values. Integers outside the range can still be converted to \gocode{Verdict}, and are refused at run time. \\
\gocode{Inscribe} and \gocode{Validate} receive no policy & Compiler & The signatures have no policy parameter. That nothing else influences them is a matter of reading and testing the code. \\
Adding a policy to \gocode{AllOf} cannot admit more records & Proof and tests & Holds for the code as written, including the fail-closed clause. \\
A change to an earlier ledger record is detected & Library under assumptions & Requires a collision-resistant hash. Rewriting the last record, truncating the file, or rewriting the whole chain is not detected without an externally held head digest. \\
A file is unchanged between discovery and verification & Not guaranteed & The size check detects some changes only. \\
Concurrent goroutines in one process cannot fork the ledger & Proof and tests & Requires that one \gocode{Writer} is the only appender to the file. Two processes can still fork the chain, as the experiment of Chapter~\ref{ch:writer} shows. \\
The pool returns one outcome per file, in input order & Proof and tests & Cancellation stops new work only; a verification in progress is not interrupted. \\
The code is free of data races & Detector and argument & The race detector sees only the executions that were run. \\
A reader of a replaced file sees the whole old or whole new contents & Operating system, with tests & Relies on rename within one filesystem. A crash can leave a stray temporary file, and durability of the rename depends on the directory flush. \\
A record that was written can be read back & Convention, with a check & \gocode{Append} refuses records above \gocode{MaxLine}; the limit is not enforced by the type system. \\
Rebuilding the index gives the same bytes as updating it & Proof and tests & Holds while \gocode{Apply} stays free of clocks, files, and map iteration order. \\
The index file has not been edited & Not guaranteed & Only its position in the ledger is checked; entries are trusted. Decisions never depend on it alone. \\
Routing depends only on the method and the path & Compiler, with reading & \gocode{Match} takes only those two strings. That nothing else enters its decision is read from thirty lines. \\
A handler runs only for a caller that \gocode{Permit} allowed & Proof and tests & Local to the package: it holds while \gocode{ServeHTTP} is the only way in. An undeclared requirement is refused. \\
The ledger has a single writer when the service is used & Convention & Holds inside one process, and only if no command-line process or second container appends to the same file. \\
A bearer token is confidential in transit & Not guaranteed & The service speaks plain HTTP and must be placed behind transport encryption or on a loopback interface. \\
A container decides what the service may do & Not guaranteed & It limits what the process can see and consume. Permission is decided by the adjudication in the program. \\
A failed append leaves no fragment in the ledger & Tests, with a real full-disk run & Holds when the truncation succeeds. A crash between the write and the cut leaves a torn final line, which the next load reports. \\
An error from \gocode{Append} means the record is absent & Not guaranteed & After a failed flush the record may be present. Callers that retry can record twice. \\
No credential reaches the log, and no file path reaches a client & Tests & The log fields are chosen one by one; a new field must be reviewed. \\
Two runs with equal manifests used the same build & Not guaranteed & The manifest holds what the toolchain records. Options such as symbol stripping, and unmodified source outside version control, are not distinguished. \\
A summary that replaces the ledger answers the decisions of the future & Not guaranteed & Shown false by exact counterexamples for the per-path index. A summary kept beside the ledger can always be recomputed. \\
\bottomrule
\end{tabular}
\caption{What is guaranteed by what, as of the end of Part~VI.}
\label{tab:guarantees}
\end{table}


\end{document}
